% Môn: Vật lý
% Lớp: 12
% Chương: Chương I. Vật lí nhiệt
% Bài: L12C1 Bài 3. Nhiệt độ. Thang nhiệt độ – nhiệt kế
% Số câu: 107
% App đọc file này trực tiếp — sửa rồi Commit trên GitHub, bấm Đồng bộ GitHub .tex

% L12C1 Bài 3. Nhiệt độ. Thang nhiệt độ – nhiệt kế

% ===== Câu 1 =====
% ID: L12C1B3-03-DS
% Mức: TH
\dangbt{Chuyển đổi nhiệt độ giữa các thang đo thông dụng}
\begin{ex}
% ID: L12C1B3-03-DS
% Mức: TH
Xét các phát biểu sau về dạng “Bài toán chuyển đổi giữa ba thang Celsius, Kelvin và Fahrenheit” (bộ số 5).
\choiceTF
{Độ tăng 1 $K$ bằng độ tăng 1 °F.}
{\True Khi giải dạng “Bài toán chuyển đổi giữa ba thang Celsius, Kelvin và Fahrenheit”, cần xác định đúng trạng thái đầu, trạng thái cuối và quy ước dấu hoặc điều kiện áp dụng.}
{Có thể bỏ qua mọi dữ kiện về khối lượng, nhiệt độ và hiệu suất trong mọi bài toán thuộc dạng này.}
{Độ tăng 10 ° $C$ bằng độ tăng 10 $K$ và 18 °F.}
\loigiai{
\begin{itemchoice}
\item (Sai) Độ tăng 1 $K$ bằng độ tăng 1 °F. (Vì): Độ tăng 1 $K$ tương đương với độ tăng 1 ° $C$, còn độ tăng 1 ° $F$ tương đương với độ tăng $\frac{5}{9}K$
\item (Đúng) Khi giải dạng “Bài toán chuyển đổi giữa ba thang Celsius, Kelvin và Fahrenheit”, cần xác định đúng trạng thái đầu, trạng thái cuối và quy ước dấu hoặc điều kiện áp dụng. (Vì): Khi giải các bài toán chuyển đổi nhiệt độ, việc xác định chính xác trạng thái ban đầu, trạng thái cuối, cũng như quy ước về dấu hoặc điều kiện áp dụng là rất quan trọng để có kết quả đúng
\item (Sai) Có thể bỏ qua mọi dữ kiện về khối lượng, nhiệt độ và hiệu suất trong mọi bài toán thuộc dạng này. (Vì): Các dữ kiện về khối lượng, nhiệt độ và hiệu suất có thể cần thiết trong một số bài toán vật lý, không thể bỏ qua tùy tiện trong mọi trường hợp
\item (Sai) Độ tăng 10 ° $C$ bằng độ tăng 10 $K$ và 18 °F.
\end{itemchoice}
}
\end{ex}

% ===== Câu 2 =====
% ID: L12C1B3-11-DS
% Mức: VD
\begin{ex}
% ID: L12C1B3-11-DS
% Mức: VD
Xét các phát biểu sau về dạng “Chuyển đổi nhiệt độ giữa thang Celsius và Fahrenheit” (bộ số 2).
\choiceTF
{\True Khi giải dạng “Chuyển đổi nhiệt độ giữa thang Celsius và Fahrenheit”, cần xác định đúng trạng thái đầu, trạng thái cuối và quy ước dấu hoặc điều kiện áp dụng.}
{Có thể bỏ qua mọi dữ kiện về khối lượng, nhiệt độ và hiệu suất trong mọi bài toán thuộc dạng này.}
{\True $t_F=1,8t_C+32.$}
{100 ° $C$ tương ứng 100 °F.}
\loigiai{
\begin{itemchoice}
\item (Đúng) Khi giải dạng “Chuyển đổi nhiệt độ giữa thang Celsius và Fahrenheit”, cần xác định đúng trạng thái đầu, trạng thái cuối và quy ước dấu hoặc điều kiện áp dụng. (Vì): Khi giải dạng “Chuyển đổi nhiệt độ giữa thang Celsius và Fahrenheit”, cần xác định đúng trạng thái đầu, trạng thái cuối và quy ước dấu hoặc điều kiện áp dụng. b) S: Có thể bỏ qua mọi dữ kiện về khối lượng, nhiệt độ và hiệu suất trong mọi bài toán thuộc dạng này. c) Đ: t_F=1,8t_C+32. d) S: 100 ° $C$ tương ứng 100 °F. Kết luận dựa trên định nghĩa, điều kiện áp dụng và quy ước trong SGK KNTT
\item (Sai) Có thể bỏ qua mọi dữ kiện về khối lượng, nhiệt độ và hiệu suất trong mọi bài toán thuộc dạng này.
\item (Đúng) $t_F=1,8t_C+32.$.
\item (Sai) 100 ° $C$ tương ứng 100 °F.
\end{itemchoice}
}
\end{ex}

% ===== Câu 3 =====
% ID: L12C1B3-12-DS
% Mức: VDC
\begin{ex}
% ID: L12C1B3-12-DS
% Mức: VDC
Xét các phát biểu sau về dạng “Chuyển đổi nhiệt độ giữa thang Celsius và Fahrenheit”
\choiceTF
{Có thể bỏ qua mọi dữ kiện về khối lượng, nhiệt độ và hiệu suất trong mọi bài toán thuộc dạng này.}
{\True $t_F=1,8t_C+32.$}
{100 ° $C$ tương ứng 100 °F.}
{Khi giải dạng “Chuyển đổi nhiệt độ giữa thang Celsius và Fahrenheit”, cần xác định đúng trạng thái đầu, trạng thái cuối và quy ước dấu hoặc điều kiện áp dụng.}
\loigiai{
\begin{itemchoice}
\item (Sai) Có thể bỏ qua mọi dữ kiện về khối lượng, nhiệt độ và hiệu suất trong mọi bài toán thuộc dạng này. (Vì): Khối lượng, nhiệt độ và hiệu suất là những yếu tố quan trọng trong nhiều bài toán vật lý, không thể bỏ qua tùy tiện. Trong dạng bài chuyển đổi nhiệt độ, nhiệt độ là dữ kiện cốt lõi, còn khối lượng và hiệu suất có thể liên quan đến các bài toán nhiệt động lực học phức tạp hơn
\item (Đúng) $t_F=1,8t_C+32.$. (Vì): Công thức chuyển đổi từ độ Celsius ($t_C$) sang độ Fahrenheit ($t_F$) là $t_F = 1,8 t_C + 32$
\item (Sai) 100 ° $C$ tương ứng 100 °F. (Vì): 100 ° $C$ tương ứng với 212 ° $F$ (nhiệt độ sôi của nước ở áp suất chuẩn), không phải 100 °F
\item (Sai) Khi giải dạng “Chuyển đổi nhiệt độ giữa thang Celsius và Fahrenheit”, cần xác định đúng trạng thái đầu, trạng thái cuối và quy ước dấu hoặc điều kiện áp dụng.
\end{itemchoice}
}
\end{ex}

% ===== Câu 4 =====
% ID: L12C1B3-13-DS
% Mức: NB
\begin{ex}
% ID: L12C1B3-13-DS
% Mức: NB
Xét các phát biểu sau về dạng “Chuyển đổi nhiệt độ giữa thang Celsius và Kelvin” (bộ số 4).
\choiceTF
{\True $T(K)=t(°C)+273,15.$}
{0 ° $C$ tương ứng 0 K.}
{\True Khi giải dạng “Chuyển đổi nhiệt độ giữa thang Celsius và Kelvin”, cần xác định đúng trạng thái đầu, trạng thái cuối và quy ước dấu hoặc điều kiện áp dụng.}
{Có thể bỏ qua mọi dữ kiện về khối lượng, nhiệt độ và hiệu suất trong mọi bài toán thuộc dạng này.}
\loigiai{
\begin{itemchoice}
\item (Đúng) $T(K)=t(°C)+273,15.$. (Vì): T(K)=t(°C)+273,15. b) S: 0 ° $C$ tương ứng 0 K. c) Đ: Khi giải dạng “Chuyển đổi nhiệt độ giữa thang Celsius và Kelvin”, cần xác định đúng trạng thái đầu, trạng thái cuối và quy ước dấu hoặc điều kiện áp dụng. d) S: Có thể bỏ qua mọi dữ kiện về khối lượng, nhiệt độ và hiệu suất trong mọi bài toán thuộc dạng này. Kết luận dựa trên định nghĩa, điều kiện áp dụng và quy ước trong SGK KNTT
\item (Sai) 0 ° $C$ tương ứng 0 K.
\item (Đúng) Khi giải dạng “Chuyển đổi nhiệt độ giữa thang Celsius và Kelvin”, cần xác định đúng trạng thái đầu, trạng thái cuối và quy ước dấu hoặc điều kiện áp dụng.
\item (Sai) Có thể bỏ qua mọi dữ kiện về khối lượng, nhiệt độ và hiệu suất trong mọi bài toán thuộc dạng này.
\end{itemchoice}
}
\end{ex}

% ===== Câu 5 =====
% ID: L12C1B3-13-TN
% Mức: NB
\begin{ex}
% ID: L12C1B3-13-TN
% Mức: NB
Chỉ số 0 nhiệt độ của một vật khi ở trạng thái cân bằng nhiệt theo thang nhiệt độ Celsius so với nhiệt độ của vật đó tính theo thang nhiệt độ Kelvin sẽ
\choice
{thấp hơn chính xác là 273,15 độ}
{\True cao hơn chính xác là 273,15 độ}
{thấp hơn chính xác là 273,16 độ}
{cao hơn chính xác là 273,16 độ}
\loigiai{
- Mối liên hệ giữa thang nhiệt độ Celsius ($t$) và Kelvin ($T$) là $T = t(^{\circ}C) + 273.15$.
 - Khi nhiệt độ là $0^{\circ}C$, nhiệt độ tương ứng theo thang Kelvin là $T = 0 + 273.15 = 273.15 K$.
 - Do đó, nhiệt độ tính theo thang Kelvin cao hơn nhiệt độ tính theo thang Celsius chính xác là 273.15 đơn vị.
}
\end{ex}

% ===== Câu 6 =====
% ID: L12C1B3-14-DS
% Mức: TH
\begin{ex}
% ID: L12C1B3-14-DS
% Mức: TH
Xét các phát biểu sau về dạng “Chuyển đổi nhiệt độ giữa thang Celsius và Kelvin” (bộ số 1).
\choiceTF
{0 ° $C$ tương ứng 0 K.}
{\True Khi giải dạng “Chuyển đổi nhiệt độ giữa thang Celsius và Kelvin”, cần xác định đúng trạng thái đầu, trạng thái cuối và quy ước dấu hoặc điều kiện áp dụng.}
{Có thể bỏ qua mọi dữ kiện về khối lượng, nhiệt độ và hiệu suất trong mọi bài toán thuộc dạng này.}
{\True $T(K)=t(°C)+273,15.$}
\loigiai{
\begin{itemchoice}
\item (Sai) 0 ° $C$ tương ứng 0 K. (Vì): 0 ° $C$ tương ứng 0 K. b) Đ: Khi giải dạng “Chuyển đổi nhiệt độ giữa thang Celsius và Kelvin”, cần xác định đúng trạng thái đầu, trạng thái cuối và quy ước dấu hoặc điều kiện áp dụng. c) S: Có thể bỏ qua mọi dữ kiện về khối lượng, nhiệt độ và hiệu suất trong mọi bài toán thuộc dạng này. d) Đ: T(K)=t(°C)+273,15. Kết luận dựa trên định nghĩa, điều kiện áp dụng và quy ước trong SGK KNTT
\item (Đúng) Khi giải dạng “Chuyển đổi nhiệt độ giữa thang Celsius và Kelvin”, cần xác định đúng trạng thái đầu, trạng thái cuối và quy ước dấu hoặc điều kiện áp dụng.
\item (Sai) Có thể bỏ qua mọi dữ kiện về khối lượng, nhiệt độ và hiệu suất trong mọi bài toán thuộc dạng này.
\item (Đúng) $T(K)=t(°C)+273,15.$.
\end{itemchoice}
}
\end{ex}

% ===== Câu 7 =====
% ID: L12C1B3-15-DS
% Mức: TH
\begin{ex}
% ID: L12C1B3-15-DS
% Mức: TH
Xét các phát biểu sau về dạng “Chuyển đổi nhiệt độ giữa thang Celsius và Kelvin” (bộ số 5).
\choiceTF
{0 ° $C$ tương ứng 0 K.}
{\True Khi giải dạng “Chuyển đổi nhiệt độ giữa thang Celsius và Kelvin”, cần xác định đúng trạng thái đầu, trạng thái cuối và quy ước dấu hoặc điều kiện áp dụng.}
{Có thể bỏ qua mọi dữ kiện về khối lượng, nhiệt độ và hiệu suất trong mọi bài toán thuộc dạng này.}
{\True $T(K)=t(°C)+273,15.$}
\loigiai{
\begin{itemchoice}
\item (Sai) 0 ° $C$ tương ứng 0 K. (Vì): 0 ° $C$ tương ứng 0 K. b) Đ: Khi giải dạng “Chuyển đổi nhiệt độ giữa thang Celsius và Kelvin”, cần xác định đúng trạng thái đầu, trạng thái cuối và quy ước dấu hoặc điều kiện áp dụng. c) S: Có thể bỏ qua mọi dữ kiện về khối lượng, nhiệt độ và hiệu suất trong mọi bài toán thuộc dạng này. d) Đ: T(K)=t(°C)+273,15. Kết luận dựa trên định nghĩa, điều kiện áp dụng và quy ước trong SGK KNTT
\item (Đúng) Khi giải dạng “Chuyển đổi nhiệt độ giữa thang Celsius và Kelvin”, cần xác định đúng trạng thái đầu, trạng thái cuối và quy ước dấu hoặc điều kiện áp dụng.
\item (Sai) Có thể bỏ qua mọi dữ kiện về khối lượng, nhiệt độ và hiệu suất trong mọi bài toán thuộc dạng này.
\item (Đúng) $T(K)=t(°C)+273,15.$.
\end{itemchoice}
}
\end{ex}

% ===== Câu 8 =====
% ID: L12C1B3-18-TL
% Mức: VDC
\begin{bt}
% ID: L12C1B3-18-TL
% Mức: VDC
Trình bày cách nhận biết và phương pháp giải dạng “Chuyển đổi nhiệt độ giữa thang Celsius và Kelvin”. Phân tích nhận định sau và sửa lại nếu sai: “T(K)=t(°C)+273,15.”
\loigiai{
Trước hết xác định hiện tượng, trạng thái và các đại lượng đã cho. Nêu định nghĩa hoặc hệ thức phù hợp, đổi về đơn vị SI, áp dụng đúng điều kiện và kiểm tra ý nghĩa kết quả. Nhận định cần đối chiếu: T(K)=t(°C)+273,15. Vì vậy nhận định đã cho đúng.
}
\end{bt}

% ===== Câu 9 =====
% ID: L12C1B3-20-TN
% Mức: VD
\dangbt{Thiết lập và chuyển đổi thang nhiệt độ tự chế}
\begin{ex}
% ID: L12C1B3-20-TN
% Mức: VD
Dùng một nhiệt kế đo theo thang nhiệt độ $Z$ tự chế (có mốc nước đá đang tan là -5 độ $Z$ và nước sôi là 105 độ Z) để đo nhiệt độ của một vật thì thấy số chỉ là 61 độ Z. Hỏi nhiệt độ của vật này trong thang Celsius là bao nhiêu?
\choice
{50 độ $C$}
{\True 60 độ $C$}
{55 độ $C$}
{65 độ $C$}
\loigiai{
Từ công thức liên hệ đã thiết lập giữa thang Celsius và thang Z:
$T = 1,1 \cdot t - 5$

Với $T = 61$ độ $Z$, ta thay vào phương trình để tìm $t$ (độ C):

$61 = 1,1 \cdot t - 5\Leftrightarrow 1,1 \cdot t = 66\Leftrightarrow t = 60^\circ\mathrm{C}$

Vậy nhiệt độ của vật tương ứng là $60^\circ\mathrm{C}$.
}
\end{ex}

% ===== Câu 10 =====
% ID: L12C1B3-21-DS
% Mức: VD
\dangbt{Lý thuyết về nhiệt độ và cân bằng nhiệt}
\begin{ex}
% ID: L12C1B3-21-DS
% Mức: VD
Xét các phát biểu sau về dạng “Lí thuyết về nhiệt độ, cân bằng nhiệt và độ không tuyệt đối” (bộ số 2).
\choiceTF
{\True Khi giải dạng “Lí thuyết về nhiệt độ, cân bằng nhiệt và độ không tuyệt đối”, cần xác định đúng trạng thái đầu, trạng thái cuối và quy ước dấu hoặc điều kiện áp dụng.}
{Có thể bỏ qua mọi dữ kiện về khối lượng, nhiệt độ và hiệu suất trong mọi bài toán thuộc dạng này.}
{\True Hai vật cân bằng nhiệt có cùng nhiệt độ.}
{Có thể đạt nhiệt độ thấp hơn 0 K.}
\loigiai{
\begin{itemchoice}
\item (Đúng) Khi giải dạng “Lí thuyết về nhiệt độ, cân bằng nhiệt và độ không tuyệt đối”, cần xác định đúng trạng thái đầu, trạng thái cuối và quy ước dấu hoặc điều kiện áp dụng. (Vì): Khi giải dạng “Lí thuyết về nhiệt độ, cân bằng nhiệt và độ không tuyệt đối”, cần xác định đúng trạng thái đầu, trạng thái cuối và quy ước dấu hoặc điều kiện áp dụng. b) S: Có thể bỏ qua mọi dữ kiện về khối lượng, nhiệt độ và hiệu suất trong mọi bài toán thuộc dạng này. c) Đ: Hai vật cân bằng nhiệt có cùng nhiệt độ. d) S: Có thể đạt nhiệt độ thấp hơn 0 K. Kết luận dựa trên định nghĩa, điều kiện áp dụng và quy ước trong SGK KNTT
\item (Sai) Có thể bỏ qua mọi dữ kiện về khối lượng, nhiệt độ và hiệu suất trong mọi bài toán thuộc dạng này.
\item (Đúng) Hai vật cân bằng nhiệt có cùng nhiệt độ.
\item (Sai) Có thể đạt nhiệt độ thấp hơn 0 K.
\end{itemchoice}
}
\end{ex}

% ===== Câu 11 =====
% ID: L12C1B3-21-TL
% Mức: TH
\begin{bt}
% ID: L12C1B3-21-TL
% Mức: TH
Trình bày cách nhận biết và phương pháp giải dạng “Lí thuyết về nhiệt độ, cân bằng nhiệt và độ không tuyệt đối”. Phân tích nhận định sau và sửa lại nếu sai: “Hai vật cân bằng nhiệt có cùng nhiệt độ.”
\loigiai{
Trước hết xác định hiện tượng, trạng thái và các đại lượng đã cho. Nêu định nghĩa hoặc hệ thức phù hợp, đổi về đơn vị SI, áp dụng đúng điều kiện và kiểm tra ý nghĩa kết quả. Nhận định cần đối chiếu: Hai vật cân bằng nhiệt có cùng nhiệt độ. Vì vậy nhận định đã cho đúng.
}
\end{bt}

% ===== Câu 12 =====
% ID: L12C1B3-21-TLN
% Mức: VDC
\dangbt{Chuyển đổi nhiệt độ giữa các thang đo thông dụng}
\begin{ex}
% ID: L12C1B3-21-TLN
% Mức: VDC
Cho bốn nhận định: (1) T(K)=t(°C)+273,15. (2) 0 ° $C$ tương ứng 0 K. (3) Cần kiểm tra đơn vị và điều kiện áp dụng khi xử lí dạng “Chuyển đổi nhiệt độ giữa thang Celsius và Kelvin”. (4) Có thể thay tùy ý nhiệt độ Celsius cho nhiệt độ tuyệt đối trong mọi công thức. Có bao nhiêu nhận định đúng? Trả lời bằng một số nguyên.
\shortans{2}
\loigiai{
Nhận định (1) và (3) đúng; (2) và (4) sai. Vậy có 2 nhận định đúng.
}
\end{ex}

% ===== Câu 13 =====
% ID: L12C1B3-21-TN
% Mức: VDC
\dangbt{Thiết lập và chuyển đổi thang nhiệt độ tự chế}
\begin{ex}
% ID: L12C1B3-21-TN
% Mức: VDC
Với thang nhiệt độ $Z$ tự chế (chọn mốc nước đá đang tan ở 1 atm là -5 độ $Z$ và mốc nước sôi là 105 độ Z), ở nhiệt độ nào của vật thì số chỉ trên hai thang nhiệt độ Celsius và thang $Z$ bằng nhau?
\choice
{\True 50}
{-50}
{-25}
{100}
\loigiai{
Gọi $t$ là nhiệt độ theo thang Celsius và $T$ là nhiệt độ theo thang $Z$.
Ta có mốc nước đá tan ở $0^\circ\mathrm{C}$ ứng với $-5^\circ\mathrm{Z}$ và mốc nước sôi ở $100^\circ\mathrm{C}$ ứng với $105^\circ\mathrm{Z}$.
Mối liên hệ giữa hai thang nhiệt độ là một hàm bậc nhất: $T = at + b$.
Thay các giá trị đã biết vào phương trình:
Khi $t = 0^\circ\mathrm{C}$ thì $T = -5^\circ\mathrm{Z}$, ta có: $-5 = a(0) + b \Rightarrow b = -5$.
Khi $t = 100^\circ\mathrm{C}$ thì $T = 105^\circ\mathrm{Z}$, ta có: $105 = a(100) - 5 \Rightarrow 100a = 110 \Rightarrow a = 1,1$.
Vậy công thức liên hệ giữa hai thang nhiệt độ là $T = 1,1t - 5$.
Để số chỉ trên hai thang nhiệt độ bằng nhau, ta đặt $t = T = x$.
Thay vào công thức liên hệ: $x = 1,1x - 5$.
Giải phương trình: $1,1x - x = 5 \Rightarrow 0,1x = 5 \Rightarrow x = 50$.
Vậy ở nhiệt độ $50^\circ\mathrm{C}$ (tương đương $50$ độ Z) thì số chỉ của hai thang nhiệt độ bằng nhau.
}
\end{ex}

% ===== Câu 14 =====
% ID: L12C1B3-22-TL
% Mức: VD
\dangbt{Lý thuyết về nhiệt độ và cân bằng nhiệt}
\begin{bt}
% ID: L12C1B3-22-TL
% Mức: VD
Trình bày cách nhận biết và phương pháp giải dạng “Lí thuyết về nhiệt độ, cân bằng nhiệt và độ không tuyệt đối”. Phân tích nhận định sau và sửa lại nếu sai: “Có thể đạt nhiệt độ thấp hơn 0 K.”
\loigiai{
Trước hết xác định hiện tượng, trạng thái và các đại lượng đã cho. Nêu định nghĩa hoặc hệ thức phù hợp, đổi về đơn vị SI, áp dụng đúng điều kiện và kiểm tra ý nghĩa kết quả. Nhận định cần đối chiếu: Hai vật cân bằng nhiệt có cùng nhiệt độ. Vì vậy nhận định đã cho sai; phát biểu đúng là: Hai vật cân bằng nhiệt có cùng nhiệt độ.
}
\end{bt}

% ===== Câu 15 =====
% ID: L12C1B3-22-TLN
% Mức: NB
\begin{ex}
% ID: L12C1B3-22-TLN
% Mức: NB
Cho bốn nhận định: (1) Hai vật cân bằng nhiệt có cùng nhiệt độ. (2) Có thể đạt nhiệt độ thấp hơn 0 K. (3) Cần kiểm tra đơn vị và điều kiện áp dụng khi xử lí dạng “Lí thuyết về nhiệt độ, cân bằng nhiệt và độ không tuyệt đối”. (4) Có thể thay tùy ý nhiệt độ Celsius cho nhiệt độ tuyệt đối trong mọi công thức. Có bao nhiêu nhận định đúng? Trả lời bằng một số nguyên.
\shortans{2}
\loigiai{
Nhận định (1) và (3) đúng; (2) và (4) sai. Vậy có 2 nhận định đúng.
}
\end{ex}

% ===== Câu 16 =====
% ID: L12C1B3-22-TN
% Mức: VDC
\dangbt{Thiết lập và chuyển đổi thang nhiệt độ tự chế}
\begin{ex}
% ID: L12C1B3-22-TN
% Mức: VDC
Một học sinh chế tạo một thang nhiệt độ tự chọn kí hiệu là thang đo X. Thang đo này lấy nhiệt độ nước đá đang tan ở áp suất tiêu chuẩn là $-10^\circ\mathrm{X}$. Nhiệt độ cơ thể của một người bình thường là $37^\circ\mathrm{C}$ tương ứng với $40^\circ\mathrm{X}$ trên thang mới này. Hãy xác định nhiệt độ sôi của nước ($100^\circ\mathrm{C}$) trên thang đo X.
\choice
{115 độ $X$}
{120 độ $X$}
{\True 125 độ $X$}
{130 độ $X$}
\loigiai{
Giả sử biểu thức chuyển đổi giữa thang Celsius $t$ ($^\circ\mathrm{C}$) và thang $X$ là hàm số bậc nhất:
$T = a \cdot t + b$

- Tại mốc nước đá đang tan $t = 0^\circ\mathrm{C}$, ta có $T = -10^\circ\mathrm{X}$. Thay vào biểu thức:
$-10 = a \cdot 0 + b \Rightarrow b = -10$.

- Tại mốc thân nhiệt người $t = 37^\circ\mathrm{C}$, ta có $T = 40^\circ\mathrm{X}$. Thay vào biểu thức:
$40 = a \cdot 37 - 10 \Leftrightarrow 37a = 50 \Leftrightarrow a = \dfrac{50}{37}$.

Như vậy, công thức liên hệ là:
$T = \dfrac{50}{37} \cdot t - 10$.

Để tìm nhiệt độ sôi của nước ($t = 100^\circ\mathrm{C}$) trên thang $X$, ta thay $t = 100$ vào công thức trên:
$T = \dfrac{50}{37} \cdot 100 - 10 = \dfrac{5000}{37} - 10 \approx 125^\circ\mathrm{X}$.
}
\end{ex}

% ===== Câu 17 =====
% ID: L12C1B3-23-DS
% Mức: NB
\dangbt{Nguyên lý hoạt động và sử dụng nhiệt kế}
\begin{ex}
% ID: L12C1B3-23-DS
% Mức: NB
Xét các phát biểu sau về dạng “Nguyên lí hoạt động và sử dụng nhiệt kế” (bộ số 4).
\choiceTF
{\True Nhiệt kế hoạt động dựa trên đại lượng vật lí biến đổi xác định theo nhiệt độ.}
{Có thể dùng mọi nhiệt kế để đo mọi khoảng nhiệt độ.}
{\True Khi giải dạng “Nguyên lí hoạt động và sử dụng nhiệt kế”, cần xác định đúng trạng thái đầu, trạng thái cuối và quy ước dấu hoặc điều kiện áp dụng.}
{Có thể bỏ qua mọi dữ kiện về khối lượng, nhiệt độ và hiệu suất trong mọi bài toán thuộc dạng này.}
\loigiai{
\begin{itemchoice}
\item (Đúng) Nhiệt kế hoạt động dựa trên đại lượng vật lí biến đổi xác định theo nhiệt độ. (Vì): Nhiệt kế hoạt động dựa trên đại lượng vật lí biến đổi xác định theo nhiệt độ. b) S: Có thể dùng mọi nhiệt kế để đo mọi khoảng nhiệt độ. c) Đ: Khi giải dạng “Nguyên lí hoạt động và sử dụng nhiệt kế”, cần xác định đúng trạng thái đầu, trạng thái cuối và quy ước dấu hoặc điều kiện áp dụng. d) S: Có thể bỏ qua mọi dữ kiện về khối lượng, nhiệt độ và hiệu suất trong mọi bài toán thuộc dạng này. Kết luận dựa trên định nghĩa, điều kiện áp dụng và quy ước trong SGK KNTT
\item (Sai) Có thể dùng mọi nhiệt kế để đo mọi khoảng nhiệt độ.
\item (Đúng) Khi giải dạng “Nguyên lí hoạt động và sử dụng nhiệt kế”, cần xác định đúng trạng thái đầu, trạng thái cuối và quy ước dấu hoặc điều kiện áp dụng.
\item (Sai) Có thể bỏ qua mọi dữ kiện về khối lượng, nhiệt độ và hiệu suất trong mọi bài toán thuộc dạng này.
\end{itemchoice}
}
\end{ex}

% ===== Câu 18 =====
% ID: L12C1B3-23-TL
% Mức: VDC

% ===== Câu 20 =====
% ID: L12C1B3-23-TN
% Mức: NB
\dangbt{Chuyển đổi nhiệt độ giữa các thang đo thông dụng}
\begin{ex}
% ID: L12C1B3-23-TN
% Mức: NB
Câu 1 thuộc dạng “Chuyển đổi nhiệt độ giữa thang Celsius và Fahrenheit”. Phát biểu nào sau đây đúng?
\choice
{100 ° $C$ tương ứng 100 °F.}
{Nhiệt lượng và nhiệt độ là hai đại lượng hoàn toàn giống nhau.}
{Mọi quá trình nhiệt học đều không phụ thuộc điều kiện thực hiện.}
{\True $t_F=1,8t_C+32.$}
\loigiai{
Phát biểu đúng là: t_F=1,8t_C+32. Các phương án còn lại trái với mô hình, định nghĩa hoặc hệ thức nhiệt học tương ứng.
}
\end{ex}

% ===== Câu 21 =====
% ID: L12C1B3-24-DS
% Mức: TH
\dangbt{Nguyên lý hoạt động và sử dụng nhiệt kế}
\begin{ex}
% ID: L12C1B3-24-DS
% Mức: TH
Xét các phát biểu sau về dạng “Nguyên lí hoạt động và sử dụng nhiệt kế” (bộ số 1).
\choiceTF
{Có thể dùng mọi nhiệt kế để đo mọi khoảng nhiệt độ.}
{\True Khi giải dạng “Nguyên lí hoạt động và sử dụng nhiệt kế”, cần xác định đúng trạng thái đầu, trạng thái cuối và quy ước dấu hoặc điều kiện áp dụng.}
{Có thể bỏ qua mọi dữ kiện về khối lượng, nhiệt độ và hiệu suất trong mọi bài toán thuộc dạng này.}
{\True Nhiệt kế hoạt động dựa trên đại lượng vật lí biến đổi xác định theo nhiệt độ.}
\loigiai{
\begin{itemchoice}
\item (Sai) Có thể dùng mọi nhiệt kế để đo mọi khoảng nhiệt độ. (Vì): Có thể dùng mọi nhiệt kế để đo mọi khoảng nhiệt độ. b) Đ: Khi giải dạng “Nguyên lí hoạt động và sử dụng nhiệt kế”, cần xác định đúng trạng thái đầu, trạng thái cuối và quy ước dấu hoặc điều kiện áp dụng. c) S: Có thể bỏ qua mọi dữ kiện về khối lượng, nhiệt độ và hiệu suất trong mọi bài toán thuộc dạng này. d) Đ: Nhiệt kế hoạt động dựa trên đại lượng vật lí biến đổi xác định theo nhiệt độ. Kết luận dựa trên định nghĩa, điều kiện áp dụng và quy ước trong SGK KNTT
\item (Đúng) Khi giải dạng “Nguyên lí hoạt động và sử dụng nhiệt kế”, cần xác định đúng trạng thái đầu, trạng thái cuối và quy ước dấu hoặc điều kiện áp dụng.
\item (Sai) Có thể bỏ qua mọi dữ kiện về khối lượng, nhiệt độ và hiệu suất trong mọi bài toán thuộc dạng này.
\item (Đúng) Nhiệt kế hoạt động dựa trên đại lượng vật lí biến đổi xác định theo nhiệt độ.
\end{itemchoice}
}
\end{ex}

% ===== Câu 22 =====
% ID: L12C1B3-24-TL
% Mức: VDC

% ===== Câu 24 =====
% ID: L12C1B3-24-TN
% Mức: NB
\dangbt{Chuyển đổi nhiệt độ giữa các thang đo thông dụng}
\begin{ex}
% ID: L12C1B3-24-TN
% Mức: NB
Câu 5 thuộc dạng “Chuyển đổi nhiệt độ giữa thang Celsius và Fahrenheit”. Phát biểu nào sau đây đúng?
\choice
{100 ° $C$ tương ứng 100 °F.}
{Nhiệt lượng và nhiệt độ là hai đại lượng hoàn toàn giống nhau.}
{Mọi quá trình nhiệt học đều không phụ thuộc điều kiện thực hiện.}
{\True $t_F=1,8t_C+32.$}
\loigiai{
Phát biểu đúng là: t_F=1,8t_C+32. Các phương án còn lại trái với mô hình, định nghĩa hoặc hệ thức nhiệt học tương ứng.
}
\end{ex}

% ===== Câu 25 =====
% ID: L12C1B3-25-DS
% Mức: TH
\dangbt{Nguyên lý hoạt động và sử dụng nhiệt kế}
\begin{ex}
% ID: L12C1B3-25-DS
% Mức: TH
Xét các phát biểu sau về dạng “Nguyên lí hoạt động và sử dụng nhiệt kế” (bộ số 5).
\choiceTF
{Có thể dùng mọi nhiệt kế để đo mọi khoảng nhiệt độ.}
{\True Khi giải dạng “Nguyên lí hoạt động và sử dụng nhiệt kế”, cần xác định đúng trạng thái đầu, trạng thái cuối và quy ước dấu hoặc điều kiện áp dụng.}
{Có thể bỏ qua mọi dữ kiện về khối lượng, nhiệt độ và hiệu suất trong mọi bài toán thuộc dạng này.}
{\True Nhiệt kế hoạt động dựa trên đại lượng vật lí biến đổi xác định theo nhiệt độ.}
\loigiai{
\begin{itemchoice}
\item (Sai) Có thể dùng mọi nhiệt kế để đo mọi khoảng nhiệt độ. (Vì): Có thể dùng mọi nhiệt kế để đo mọi khoảng nhiệt độ. b) Đ: Khi giải dạng “Nguyên lí hoạt động và sử dụng nhiệt kế”, cần xác định đúng trạng thái đầu, trạng thái cuối và quy ước dấu hoặc điều kiện áp dụng. c) S: Có thể bỏ qua mọi dữ kiện về khối lượng, nhiệt độ và hiệu suất trong mọi bài toán thuộc dạng này. d) Đ: Nhiệt kế hoạt động dựa trên đại lượng vật lí biến đổi xác định theo nhiệt độ. Kết luận dựa trên định nghĩa, điều kiện áp dụng và quy ước trong SGK KNTT
\item (Đúng) Khi giải dạng “Nguyên lí hoạt động và sử dụng nhiệt kế”, cần xác định đúng trạng thái đầu, trạng thái cuối và quy ước dấu hoặc điều kiện áp dụng.
\item (Sai) Có thể bỏ qua mọi dữ kiện về khối lượng, nhiệt độ và hiệu suất trong mọi bài toán thuộc dạng này.
\item (Đúng) Nhiệt kế hoạt động dựa trên đại lượng vật lí biến đổi xác định theo nhiệt độ.
\end{itemchoice}
}
\end{ex}

% ===== Câu 26 =====
% ID: L12C1B3-25-TL
% Mức: NB
\begin{bt}
% ID: L12C1B3-25-TL
% Mức: NB
Trình bày cách nhận biết và phương pháp giải dạng “Nguyên lí hoạt động và sử dụng nhiệt kế”. Phân tích nhận định sau và sửa lại nếu sai: “Có thể dùng mọi nhiệt kế để đo mọi khoảng nhiệt độ.”
\loigiai{
Trước hết xác định hiện tượng, trạng thái và các đại lượng đã cho. Nêu định nghĩa hoặc hệ thức phù hợp, đổi về đơn vị SI, áp dụng đúng điều kiện và kiểm tra ý nghĩa kết quả. Nhận định cần đối chiếu: Nhiệt kế hoạt động dựa trên đại lượng vật lí biến đổi xác định theo nhiệt độ. Vì vậy nhận định đã cho sai; phát biểu đúng là: Nhiệt kế hoạt động dựa trên đại lượng vật lí biến đổi xác định theo nhiệt độ.
}
\end{bt}

% ===== Câu 27 =====
% ID: L12C1B3-25-TLN
% Mức: VD

% ===== Câu 28 =====
% ID: L12C1B3-25-TN
% Mức: NB
\dangbt{Chuyển đổi nhiệt độ giữa các thang đo thông dụng}
\begin{ex}
% ID: L12C1B3-25-TN
% Mức: NB
Câu 9 thuộc dạng “Chuyển đổi nhiệt độ giữa thang Celsius và Fahrenheit”. Phát biểu nào sau đây đúng?
\choice
{100 ° $C$ tương ứng 100 °F.}
{Nhiệt lượng và nhiệt độ là hai đại lượng hoàn toàn giống nhau.}
{Mọi quá trình nhiệt học đều không phụ thuộc điều kiện thực hiện.}
{\True $t_F=1,8t_C+32.$}
\loigiai{
Phát biểu đúng là: t_F=1,8t_C+32. Các phương án còn lại trái với mô hình, định nghĩa hoặc hệ thức nhiệt học tương ứng.
}
\end{ex}

% ===== Câu 29 =====
% ID: L12C1B3-26-DS
% Mức: VD
\dangbt{Nguyên lý hoạt động và sử dụng nhiệt kế}
\begin{ex}
% ID: L12C1B3-26-DS
% Mức: VD
Xét các phát biểu sau về dạng “Nguyên lí hoạt động và sử dụng nhiệt kế” (bộ số 2).
\choiceTF
{\True Khi giải dạng “Nguyên lí hoạt động và sử dụng nhiệt kế”, cần xác định đúng trạng thái đầu, trạng thái cuối và quy ước dấu hoặc điều kiện áp dụng.}
{Có thể bỏ qua mọi dữ kiện về khối lượng, nhiệt độ và hiệu suất trong mọi bài toán thuộc dạng này.}
{\True Nhiệt kế hoạt động dựa trên đại lượng vật lí biến đổi xác định theo nhiệt độ.}
{Có thể dùng mọi nhiệt kế để đo mọi khoảng nhiệt độ.}
\loigiai{
\begin{itemchoice}
\item (Đúng) Khi giải dạng “Nguyên lí hoạt động và sử dụng nhiệt kế”, cần xác định đúng trạng thái đầu, trạng thái cuối và quy ước dấu hoặc điều kiện áp dụng. (Vì): Khi giải dạng “Nguyên lí hoạt động và sử dụng nhiệt kế”, cần xác định đúng trạng thái đầu, trạng thái cuối và quy ước dấu hoặc điều kiện áp dụng. b) S: Có thể bỏ qua mọi dữ kiện về khối lượng, nhiệt độ và hiệu suất trong mọi bài toán thuộc dạng này. c) Đ: Nhiệt kế hoạt động dựa trên đại lượng vật lí biến đổi xác định theo nhiệt độ. d) S: Có thể dùng mọi nhiệt kế để đo mọi khoảng nhiệt độ. Kết luận dựa trên định nghĩa, điều kiện áp dụng và quy ước trong SGK KNTT
\item (Sai) Có thể bỏ qua mọi dữ kiện về khối lượng, nhiệt độ và hiệu suất trong mọi bài toán thuộc dạng này.
\item (Đúng) Nhiệt kế hoạt động dựa trên đại lượng vật lí biến đổi xác định theo nhiệt độ.
\item (Sai) Có thể dùng mọi nhiệt kế để đo mọi khoảng nhiệt độ.
\end{itemchoice}
}
\end{ex}

% ===== Câu 31 =====
% ID: L12C1B3-26-TLN
% Mức: VDC

% ===== Câu 32 =====
% ID: L12C1B3-26-TN
% Mức: TH
\dangbt{Chuyển đổi nhiệt độ giữa các thang đo thông dụng}
\begin{ex}
% ID: L12C1B3-26-TN
% Mức: TH
Câu 2 thuộc dạng “Chuyển đổi nhiệt độ giữa thang Celsius và Fahrenheit”. Phát biểu nào sau đây đúng?
\choice
{Nhiệt lượng và nhiệt độ là hai đại lượng hoàn toàn giống nhau.}
{Mọi quá trình nhiệt học đều không phụ thuộc điều kiện thực hiện.}
{\True $t_F=1,8t_C+32.$}
{100 ° $C$ tương ứng 100 °F.}
\loigiai{
Phát biểu đúng là: t_F=1,8t_C+32. Các phương án còn lại trái với mô hình, định nghĩa hoặc hệ thức nhiệt học tương ứng.
}
\end{ex}

% ===== Câu 33 =====
% ID: L12C1B3-27-DS
% Mức: VDC
\dangbt{Nguyên lý hoạt động và sử dụng nhiệt kế}
\begin{ex}
% ID: L12C1B3-27-DS
% Mức: VDC
Xét các phát biểu sau về dạng “Nguyên lí hoạt động và sử dụng nhiệt kế” (bộ số 3).
\choiceTF
{Có thể bỏ qua mọi dữ kiện về khối lượng, nhiệt độ và hiệu suất trong mọi bài toán thuộc dạng này.}
{\True Nhiệt kế hoạt động dựa trên đại lượng vật lí biến đổi xác định theo nhiệt độ.}
{Có thể dùng mọi nhiệt kế để đo mọi khoảng nhiệt độ.}
{\True Khi giải dạng “Nguyên lí hoạt động và sử dụng nhiệt kế”, cần xác định đúng trạng thái đầu, trạng thái cuối và quy ước dấu hoặc điều kiện áp dụng.}
\loigiai{
\begin{itemchoice}
\item (Sai) Có thể bỏ qua mọi dữ kiện về khối lượng, nhiệt độ và hiệu suất trong mọi bài toán thuộc dạng này. (Vì): Có thể bỏ qua mọi dữ kiện về khối lượng, nhiệt độ và hiệu suất trong mọi bài toán thuộc dạng này. b) Đ: Nhiệt kế hoạt động dựa trên đại lượng vật lí biến đổi xác định theo nhiệt độ. c) S: Có thể dùng mọi nhiệt kế để đo mọi khoảng nhiệt độ. d) Đ: Khi giải dạng “Nguyên lí hoạt động và sử dụng nhiệt kế”, cần xác định đúng trạng thái đầu, trạng thái cuối và quy ước dấu hoặc điều kiện áp dụng. Kết luận dựa trên định nghĩa, điều kiện áp dụng và quy ước trong SGK KNTT
\item (Đúng) Nhiệt kế hoạt động dựa trên đại lượng vật lí biến đổi xác định theo nhiệt độ.
\item (Sai) Có thể dùng mọi nhiệt kế để đo mọi khoảng nhiệt độ.
\item (Đúng) Khi giải dạng “Nguyên lí hoạt động và sử dụng nhiệt kế”, cần xác định đúng trạng thái đầu, trạng thái cuối và quy ước dấu hoặc điều kiện áp dụng.
\end{itemchoice}
}
\end{ex}

% ===== Câu 34 =====
% ID: L12C1B3-27-TL
% Mức: TH
\begin{bt}
% ID: L12C1B3-27-TL
% Mức: TH
Trình bày cách nhận biết và phương pháp giải dạng “Nguyên lí hoạt động và sử dụng nhiệt kế”. Phân tích nhận định sau và sửa lại nếu sai: “Nhiệt kế hoạt động dựa trên đại lượng vật lí biến đổi xác định theo nhiệt độ.”
\loigiai{
Trước hết xác định hiện tượng, trạng thái và các đại lượng đã cho. Nêu định nghĩa hoặc hệ thức phù hợp, đổi về đơn vị SI, áp dụng đúng điều kiện và kiểm tra ý nghĩa kết quả. Nhận định cần đối chiếu: Nhiệt kế hoạt động dựa trên đại lượng vật lí biến đổi xác định theo nhiệt độ. Vì vậy nhận định đã cho đúng.
}
\end{bt}

% ===== Câu 35 =====
% ID: L12C1B3-27-TLN
% Mức: VDC

% ===== Câu 36 =====
% ID: L12C1B3-27-TN
% Mức: TH
\dangbt{Chuyển đổi nhiệt độ giữa các thang đo thông dụng}
\begin{ex}
% ID: L12C1B3-27-TN
% Mức: TH
Câu 6 thuộc dạng “Chuyển đổi nhiệt độ giữa thang Celsius và Fahrenheit”. Phát biểu nào sau đây đúng?
\choice
{Nhiệt lượng và nhiệt độ là hai đại lượng hoàn toàn giống nhau.}
{Mọi quá trình nhiệt học đều không phụ thuộc điều kiện thực hiện.}
{\True $t_F=1,8t_C+32.$}
{100 ° $C$ tương ứng 100 °F.}
\loigiai{
Phát biểu đúng là: t_F=1,8t_C+32. Các phương án còn lại trái với mô hình, định nghĩa hoặc hệ thức nhiệt học tương ứng.
}
\end{ex}

% ===== Câu 37 =====
% ID: L12C1B3-28-DS
% Mức: NB
\dangbt{Thiết lập và chuyển đổi thang nhiệt độ tự chế}
\begin{ex}
% ID: L12C1B3-28-DS
% Mức: NB
Xét các phát biểu sau về dạng “Thiết lập và chuyển đổi thang nhiệt độ mới” (bộ số 4).
\choiceTF
{\True Thang tuyến tính được xác định từ hai mốc và phép nội suy.}
{Khoảng chia của mọi thang nhiệt độ đều bằng nhau.}
{\True Khi giải dạng “Thiết lập và chuyển đổi thang nhiệt độ mới”, cần xác định đúng trạng thái đầu, trạng thái cuối và quy ước dấu hoặc điều kiện áp dụng.}
{Có thể bỏ qua mọi dữ kiện về khối lượng, nhiệt độ và hiệu suất trong mọi bài toán thuộc dạng này.}
\loigiai{
\begin{itemchoice}
\item (Đúng) Thang tuyến tính được xác định từ hai mốc và phép nội suy. (Vì): Thang tuyến tính được xác định từ hai mốc và phép nội suy. b) S: Khoảng chia của mọi thang nhiệt độ đều bằng nhau. c) Đ: Khi giải dạng “Thiết lập và chuyển đổi thang nhiệt độ mới”, cần xác định đúng trạng thái đầu, trạng thái cuối và quy ước dấu hoặc điều kiện áp dụng. d) S: Có thể bỏ qua mọi dữ kiện về khối lượng, nhiệt độ và hiệu suất trong mọi bài toán thuộc dạng này. Kết luận dựa trên định nghĩa, điều kiện áp dụng và quy ước trong SGK KNTT
\item (Sai) Khoảng chia của mọi thang nhiệt độ đều bằng nhau.
\item (Đúng) Khi giải dạng “Thiết lập và chuyển đổi thang nhiệt độ mới”, cần xác định đúng trạng thái đầu, trạng thái cuối và quy ước dấu hoặc điều kiện áp dụng.
\item (Sai) Có thể bỏ qua mọi dữ kiện về khối lượng, nhiệt độ và hiệu suất trong mọi bài toán thuộc dạng này.
\end{itemchoice}
}
\end{ex}

% ===== Câu 38 =====
% ID: L12C1B3-28-TL
% Mức: VD

% ===== Câu 39 =====
% ID: L12C1B3-28-TLN
% Mức: NB
\dangbt{Nguyên lý hoạt động và sử dụng nhiệt kế}
\begin{ex}
% ID: L12C1B3-28-TLN
% Mức: NB
Cho bốn nhận định: (1) Nhiệt kế hoạt động dựa trên đại lượng vật lí biến đổi xác định theo nhiệt độ. (2) Có thể dùng mọi nhiệt kế để đo mọi khoảng nhiệt độ. (3) Cần kiểm tra đơn vị và điều kiện áp dụng khi xử lí dạng “Nguyên lí hoạt động và sử dụng nhiệt kế”. (4) Có thể thay tùy ý nhiệt độ Celsius cho nhiệt độ tuyệt đối trong mọi công thức. Có bao nhiêu nhận định đúng? Trả lời bằng một số nguyên.
\shortans{2}
\loigiai{
Nhận định (1) và (3) đúng; (2) và (4) sai. Vậy có 2 nhận định đúng.
}
\end{ex}

% ===== Câu 40 =====
% ID: L12C1B3-28-TN
% Mức: TH
\dangbt{Chuyển đổi nhiệt độ giữa các thang đo thông dụng}
\begin{ex}
% ID: L12C1B3-28-TN
% Mức: TH
Câu 10 thuộc dạng “Chuyển đổi nhiệt độ giữa thang Celsius và Fahrenheit”. Phát biểu nào sau đây đúng?
\choice
{Nhiệt lượng và nhiệt độ là hai đại lượng hoàn toàn giống nhau.}
{Mọi quá trình nhiệt học đều không phụ thuộc điều kiện thực hiện.}
{\True $t_F=1,8t_C+32.$}
{100 ° $C$ tương ứng 100 °F.}
\loigiai{
Phát biểu đúng là: t_F=1,8t_C+32. Các phương án còn lại trái với mô hình, định nghĩa hoặc hệ thức nhiệt học tương ứng.
}
\end{ex}

% ===== Câu 41 =====
% ID: L12C1B3-29-DS
% Mức: TH
\dangbt{Thiết lập và chuyển đổi thang nhiệt độ tự chế}
\begin{ex}
% ID: L12C1B3-29-DS
% Mức: TH
Xét các phát biểu sau về dạng “Thiết lập và chuyển đổi thang nhiệt độ mới” (bộ số 1).
\choiceTF
{Khoảng chia của mọi thang nhiệt độ đều bằng nhau.}
{\True Khi giải dạng “Thiết lập và chuyển đổi thang nhiệt độ mới”, cần xác định đúng trạng thái đầu, trạng thái cuối và quy ước dấu hoặc điều kiện áp dụng.}
{Có thể bỏ qua mọi dữ kiện về khối lượng, nhiệt độ và hiệu suất trong mọi bài toán thuộc dạng này.}
{\True Thang tuyến tính được xác định từ hai mốc và phép nội suy.}
\loigiai{
\begin{itemchoice}
\item (Sai) Khoảng chia của mọi thang nhiệt độ đều bằng nhau. (Vì): Khoảng chia của mọi thang nhiệt độ đều bằng nhau. b) Đ: Khi giải dạng “Thiết lập và chuyển đổi thang nhiệt độ mới”, cần xác định đúng trạng thái đầu, trạng thái cuối và quy ước dấu hoặc điều kiện áp dụng. c) S: Có thể bỏ qua mọi dữ kiện về khối lượng, nhiệt độ và hiệu suất trong mọi bài toán thuộc dạng này. d) Đ: Thang tuyến tính được xác định từ hai mốc và phép nội suy. Kết luận dựa trên định nghĩa, điều kiện áp dụng và quy ước trong SGK KNTT
\item (Đúng) Khi giải dạng “Thiết lập và chuyển đổi thang nhiệt độ mới”, cần xác định đúng trạng thái đầu, trạng thái cuối và quy ước dấu hoặc điều kiện áp dụng.
\item (Sai) Có thể bỏ qua mọi dữ kiện về khối lượng, nhiệt độ và hiệu suất trong mọi bài toán thuộc dạng này.
\item (Đúng) Thang tuyến tính được xác định từ hai mốc và phép nội suy.
\end{itemchoice}
}
\end{ex}

% ===== Câu 42 =====
% ID: L12C1B3-29-TL
% Mức: VDC

% ===== Câu 44 =====
% ID: L12C1B3-29-TN
% Mức: VD
\dangbt{Chuyển đổi nhiệt độ giữa các thang đo thông dụng}
\begin{ex}
% ID: L12C1B3-29-TN
% Mức: VD
Câu 3 thuộc dạng “Chuyển đổi nhiệt độ giữa thang Celsius và Fahrenheit”. Phát biểu nào sau đây đúng?
\choice
{Mọi quá trình nhiệt học đều không phụ thuộc điều kiện thực hiện.}
{\True $t_F=1,8t_C+32.$}
{100 ° $C$ tương ứng 100 °F.}
{Nhiệt lượng và nhiệt độ là hai đại lượng hoàn toàn giống nhau.}
\loigiai{
Phát biểu đúng là: t_F=1,8t_C+32. Các phương án còn lại trái với mô hình, định nghĩa hoặc hệ thức nhiệt học tương ứng.
}
\end{ex}

% ===== Câu 45 =====
% ID: L12C1B3-30-DS
% Mức: TH
\dangbt{Thiết lập và chuyển đổi thang nhiệt độ tự chế}
\begin{ex}
% ID: L12C1B3-30-DS
% Mức: TH
Xét các phát biểu sau về dạng “Thiết lập và chuyển đổi thang nhiệt độ mới” (bộ số 5).
\choiceTF
{Khoảng chia của mọi thang nhiệt độ đều bằng nhau.}
{\True Khi giải dạng “Thiết lập và chuyển đổi thang nhiệt độ mới”, cần xác định đúng trạng thái đầu, trạng thái cuối và quy ước dấu hoặc điều kiện áp dụng.}
{Có thể bỏ qua mọi dữ kiện về khối lượng, nhiệt độ và hiệu suất trong mọi bài toán thuộc dạng này.}
{\True Thang tuyến tính được xác định từ hai mốc và phép nội suy.}
\loigiai{
\begin{itemchoice}
\item (Sai) Khoảng chia của mọi thang nhiệt độ đều bằng nhau. (Vì): Khoảng chia của mọi thang nhiệt độ đều bằng nhau. b) Đ: Khi giải dạng “Thiết lập và chuyển đổi thang nhiệt độ mới”, cần xác định đúng trạng thái đầu, trạng thái cuối và quy ước dấu hoặc điều kiện áp dụng. c) S: Có thể bỏ qua mọi dữ kiện về khối lượng, nhiệt độ và hiệu suất trong mọi bài toán thuộc dạng này. d) Đ: Thang tuyến tính được xác định từ hai mốc và phép nội suy. Kết luận dựa trên định nghĩa, điều kiện áp dụng và quy ước trong SGK KNTT
\item (Đúng) Khi giải dạng “Thiết lập và chuyển đổi thang nhiệt độ mới”, cần xác định đúng trạng thái đầu, trạng thái cuối và quy ước dấu hoặc điều kiện áp dụng.
\item (Sai) Có thể bỏ qua mọi dữ kiện về khối lượng, nhiệt độ và hiệu suất trong mọi bài toán thuộc dạng này.
\item (Đúng) Thang tuyến tính được xác định từ hai mốc và phép nội suy.
\end{itemchoice}
}
\end{ex}

% ===== Câu 46 =====
% ID: L12C1B3-30-TL
% Mức: VDC

% ===== Câu 48 =====
% ID: L12C1B3-30-TN
% Mức: VD
\dangbt{Chuyển đổi nhiệt độ giữa các thang đo thông dụng}
\begin{ex}
% ID: L12C1B3-30-TN
% Mức: VD
Câu 7 thuộc dạng “Chuyển đổi nhiệt độ giữa thang Celsius và Fahrenheit”. Phát biểu nào sau đây đúng?
\choice
{Mọi quá trình nhiệt học đều không phụ thuộc điều kiện thực hiện.}
{\True $t_F=1,8t_C+32.$}
{100 ° $C$ tương ứng 100 °F.}
{Nhiệt lượng và nhiệt độ là hai đại lượng hoàn toàn giống nhau.}
\loigiai{
Phát biểu đúng là: t_F=1,8t_C+32. Các phương án còn lại trái với mô hình, định nghĩa hoặc hệ thức nhiệt học tương ứng.
}
\end{ex}

% ===== Câu 49 =====
% ID: L12C1B3-31-DS
% Mức: VD
\dangbt{Thiết lập và chuyển đổi thang nhiệt độ tự chế}
\begin{ex}
% ID: L12C1B3-31-DS
% Mức: VD
Xét các phát biểu sau về dạng “Thiết lập và chuyển đổi thang nhiệt độ mới” (bộ số 2).
\choiceTF
{\True Khi giải dạng “Thiết lập và chuyển đổi thang nhiệt độ mới”, cần xác định đúng trạng thái đầu, trạng thái cuối và quy ước dấu hoặc điều kiện áp dụng.}
{Có thể bỏ qua mọi dữ kiện về khối lượng, nhiệt độ và hiệu suất trong mọi bài toán thuộc dạng này.}
{\True Thang tuyến tính được xác định từ hai mốc và phép nội suy.}
{Khoảng chia của mọi thang nhiệt độ đều bằng nhau.}
\loigiai{
\begin{itemchoice}
\item (Đúng) Khi giải dạng “Thiết lập và chuyển đổi thang nhiệt độ mới”, cần xác định đúng trạng thái đầu, trạng thái cuối và quy ước dấu hoặc điều kiện áp dụng. (Vì): Khi giải dạng “Thiết lập và chuyển đổi thang nhiệt độ mới”, cần xác định đúng trạng thái đầu, trạng thái cuối và quy ước dấu hoặc điều kiện áp dụng. b) S: Có thể bỏ qua mọi dữ kiện về khối lượng, nhiệt độ và hiệu suất trong mọi bài toán thuộc dạng này. c) Đ: Thang tuyến tính được xác định từ hai mốc và phép nội suy. d) S: Khoảng chia của mọi thang nhiệt độ đều bằng nhau. Kết luận dựa trên định nghĩa, điều kiện áp dụng và quy ước trong SGK KNTT
\item (Sai) Có thể bỏ qua mọi dữ kiện về khối lượng, nhiệt độ và hiệu suất trong mọi bài toán thuộc dạng này.
\item (Đúng) Thang tuyến tính được xác định từ hai mốc và phép nội suy.
\item (Sai) Khoảng chia của mọi thang nhiệt độ đều bằng nhau.
\end{itemchoice}
}
\end{ex}

% ===== Câu 50 =====
% ID: L12C1B3-31-TL
% Mức: NB
\begin{bt}
% ID: L12C1B3-31-TL
% Mức: NB
Trình bày cách nhận biết và phương pháp giải dạng “Thiết lập và chuyển đổi thang nhiệt độ mới”. Phân tích nhận định sau và sửa lại nếu sai: “Khoảng chia của mọi thang nhiệt độ đều bằng nhau.”
\loigiai{
Trước hết xác định hiện tượng, trạng thái và các đại lượng đã cho. Nêu định nghĩa hoặc hệ thức phù hợp, đổi về đơn vị SI, áp dụng đúng điều kiện và kiểm tra ý nghĩa kết quả. Nhận định cần đối chiếu: Thang tuyến tính được xác định từ hai mốc và phép nội suy. Vì vậy nhận định đã cho sai; phát biểu đúng là: Thang tuyến tính được xác định từ hai mốc và phép nội suy.
}
\end{bt}

% ===== Câu 51 =====
% ID: L12C1B3-31-TLN
% Mức: VD

% ===== Câu 52 =====
% ID: L12C1B3-31-TN
% Mức: VDC
\dangbt{Chuyển đổi nhiệt độ giữa các thang đo thông dụng}
\begin{ex}
% ID: L12C1B3-31-TN
% Mức: VDC
Câu 4 thuộc dạng “Chuyển đổi nhiệt độ giữa thang Celsius và Fahrenheit”. Phát biểu nào sau đây đúng?
\choice
{\True $t_F=1,8t_C+32.$}
{100 ° $C$ tương ứng 100 °F.}
{Nhiệt lượng và nhiệt độ là hai đại lượng hoàn toàn giống nhau.}
{Mọi quá trình nhiệt học đều không phụ thuộc điều kiện thực hiện.}
\loigiai{
Phát biểu đúng là: t_F=1,8t_C+32. Các phương án còn lại trái với mô hình, định nghĩa hoặc hệ thức nhiệt học tương ứng.
}
\end{ex}

% ===== Câu 53 =====
% ID: L12C1B3-32-DS
% Mức: VDC
\dangbt{Thiết lập và chuyển đổi thang nhiệt độ tự chế}
\begin{ex}
% ID: L12C1B3-32-DS
% Mức: VDC
Xét các phát biểu sau về dạng “Thiết lập và chuyển đổi thang nhiệt độ mới” (bộ số 3).
\choiceTF
{Có thể bỏ qua mọi dữ kiện về khối lượng, nhiệt độ và hiệu suất trong mọi bài toán thuộc dạng này.}
{\True Thang tuyến tính được xác định từ hai mốc và phép nội suy.}
{Khoảng chia của mọi thang nhiệt độ đều bằng nhau.}
{\True Khi giải dạng “Thiết lập và chuyển đổi thang nhiệt độ mới”, cần xác định đúng trạng thái đầu, trạng thái cuối và quy ước dấu hoặc điều kiện áp dụng.}
\loigiai{
\begin{itemchoice}
\item (Sai) Có thể bỏ qua mọi dữ kiện về khối lượng, nhiệt độ và hiệu suất trong mọi bài toán thuộc dạng này. (Vì): Có thể bỏ qua mọi dữ kiện về khối lượng, nhiệt độ và hiệu suất trong mọi bài toán thuộc dạng này. b) Đ: Thang tuyến tính được xác định từ hai mốc và phép nội suy. c) S: Khoảng chia của mọi thang nhiệt độ đều bằng nhau. d) Đ: Khi giải dạng “Thiết lập và chuyển đổi thang nhiệt độ mới”, cần xác định đúng trạng thái đầu, trạng thái cuối và quy ước dấu hoặc điều kiện áp dụng. Kết luận dựa trên định nghĩa, điều kiện áp dụng và quy ước trong SGK KNTT
\item (Đúng) Thang tuyến tính được xác định từ hai mốc và phép nội suy.
\item (Sai) Khoảng chia của mọi thang nhiệt độ đều bằng nhau.
\item (Đúng) Khi giải dạng “Thiết lập và chuyển đổi thang nhiệt độ mới”, cần xác định đúng trạng thái đầu, trạng thái cuối và quy ước dấu hoặc điều kiện áp dụng.
\end{itemchoice}
}
\end{ex}

% ===== Câu 55 =====
% ID: L12C1B3-32-TLN
% Mức: VDC

% ===== Câu 56 =====
% ID: L12C1B3-32-TN
% Mức: VDC
\dangbt{Chuyển đổi nhiệt độ giữa các thang đo thông dụng}
\begin{ex}
% ID: L12C1B3-32-TN
% Mức: VDC
Câu 8 thuộc dạng “Chuyển đổi nhiệt độ giữa thang Celsius và Fahrenheit”. Phát biểu nào sau đây đúng?
\choice
{\True $t_F=1,8t_C+32.$}
{100 ° $C$ tương ứng 100 °F.}
{Nhiệt lượng và nhiệt độ là hai đại lượng hoàn toàn giống nhau.}
{Mọi quá trình nhiệt học đều không phụ thuộc điều kiện thực hiện.}
\loigiai{
Phát biểu đúng là: t_F=1,8t_C+32. Các phương án còn lại trái với mô hình, định nghĩa hoặc hệ thức nhiệt học tương ứng.
}
\end{ex}

% ===== Câu 57 =====
% ID: L12C1B3-33-TL
% Mức: TH
\dangbt{Thiết lập và chuyển đổi thang nhiệt độ tự chế}
\begin{bt}
% ID: L12C1B3-33-TL
% Mức: TH
Trình bày cách nhận biết và phương pháp giải dạng “Thiết lập và chuyển đổi thang nhiệt độ mới”. Phân tích nhận định sau và sửa lại nếu sai: “Thang tuyến tính được xác định từ hai mốc và phép nội suy.”
\loigiai{
Trước hết xác định hiện tượng, trạng thái và các đại lượng đã cho. Nêu định nghĩa hoặc hệ thức phù hợp, đổi về đơn vị SI, áp dụng đúng điều kiện và kiểm tra ý nghĩa kết quả. Nhận định cần đối chiếu: Thang tuyến tính được xác định từ hai mốc và phép nội suy. Vì vậy nhận định đã cho đúng.
}
\end{bt}

% ===== Câu 58 =====
% ID: L12C1B3-33-TLN
% Mức: VDC

% ===== Câu 59 =====
% ID: L12C1B3-33-TN
% Mức: NB
\dangbt{Chuyển đổi nhiệt độ giữa các thang đo thông dụng}
\begin{ex}
% ID: L12C1B3-33-TN
% Mức: NB
Câu 1 thuộc dạng “Chuyển đổi nhiệt độ giữa thang Celsius và Kelvin”. Phát biểu nào sau đây đúng?
\choice
{0 ° $C$ tương ứng 0 K.}
{Nhiệt lượng và nhiệt độ là hai đại lượng hoàn toàn giống nhau.}
{Mọi quá trình nhiệt học đều không phụ thuộc điều kiện thực hiện.}
{\True $T(K)=t(°C)+273,15.$}
\loigiai{
Phát biểu đúng là: T(K)=t(°C)+273,15. Các phương án còn lại trái với mô hình, định nghĩa hoặc hệ thức nhiệt học tương ứng.
}
\end{ex}

% ===== Câu 60 =====
% ID: L12C1B3-34-TL
% Mức: VD

% ===== Câu 61 =====
% ID: L12C1B3-34-TLN
% Mức: NB
\dangbt{Thiết lập và chuyển đổi thang nhiệt độ tự chế}
\begin{ex}
% ID: L12C1B3-34-TLN
% Mức: NB
Cho bốn nhận định: (1) Thang tuyến tính được xác định từ hai mốc và phép nội suy. (2) Khoảng chia của mọi thang nhiệt độ đều bằng nhau. (3) Cần kiểm tra đơn vị và điều kiện áp dụng khi xử lí dạng “Thiết lập và chuyển đổi thang nhiệt độ mới”. (4) Có thể thay tùy ý nhiệt độ Celsius cho nhiệt độ tuyệt đối trong mọi công thức. Có bao nhiêu nhận định đúng? Trả lời bằng một số nguyên.
\shortans{2}
\loigiai{
Nhận định (1) và (3) đúng; (2) và (4) sai. Vậy có 2 nhận định đúng.
}
\end{ex}

% ===== Câu 62 =====
% ID: L12C1B3-34-TN
% Mức: NB
\dangbt{Chuyển đổi nhiệt độ giữa các thang đo thông dụng}
\begin{ex}
% ID: L12C1B3-34-TN
% Mức: NB
Câu 5 thuộc dạng “Chuyển đổi nhiệt độ giữa thang Celsius và Kelvin”. Phát biểu nào sau đây đúng?
\choice
{0 ° $C$ tương ứng 0 K.}
{Nhiệt lượng và nhiệt độ là hai đại lượng hoàn toàn giống nhau.}
{Mọi quá trình nhiệt học đều không phụ thuộc điều kiện thực hiện.}
{\True $T(K)=t(°C)+273,15.$}
\loigiai{
Phát biểu đúng là: T(K)=t(°C)+273,15. Các phương án còn lại trái với mô hình, định nghĩa hoặc hệ thức nhiệt học tương ứng.
}
\end{ex}

% ===== Câu 63 =====
% ID: L12C1B3-35-TL
% Mức: VDC

% ===== Câu 65 =====
% ID: L12C1B3-35-TN
% Mức: NB
\begin{ex}
% ID: L12C1B3-35-TN
% Mức: NB
Câu 9 thuộc dạng “Chuyển đổi nhiệt độ giữa thang Celsius và Kelvin”. Phát biểu nào sau đây đúng?
\choice
{0 ° $C$ tương ứng 0 K.}
{Nhiệt lượng và nhiệt độ là hai đại lượng hoàn toàn giống nhau.}
{Mọi quá trình nhiệt học đều không phụ thuộc điều kiện thực hiện.}
{\True $T(K)=t(°C)+273,15.$}
\loigiai{
Phát biểu đúng là: T(K)=t(°C)+273,15. Các phương án còn lại trái với mô hình, định nghĩa hoặc hệ thức nhiệt học tương ứng.
}
\end{ex}

% ===== Câu 66 =====
% ID: L12C1B3-36-TL
% Mức: VDC

% ===== Câu 68 =====
% ID: L12C1B3-36-TN
% Mức: TH
\begin{ex}
% ID: L12C1B3-36-TN
% Mức: TH
Câu 2 thuộc dạng “Chuyển đổi nhiệt độ giữa thang Celsius và Kelvin”. Phát biểu nào sau đây đúng?
\choice
{Nhiệt lượng và nhiệt độ là hai đại lượng hoàn toàn giống nhau.}
{Mọi quá trình nhiệt học đều không phụ thuộc điều kiện thực hiện.}
{\True $T(K)=t(°C)+273,15.$}
{0 ° $C$ tương ứng 0 K.}
\loigiai{
Phát biểu đúng là: T(K)=t(°C)+273,15. Các phương án còn lại trái với mô hình, định nghĩa hoặc hệ thức nhiệt học tương ứng.
}
\end{ex}

% ===== Câu 69 =====
% ID: L12C1B3-37-TLN
% Mức: VD

% ===== Câu 70 =====
% ID: L12C1B3-37-TN
% Mức: TH
\begin{ex}
% ID: L12C1B3-37-TN
% Mức: TH
Câu 6 thuộc dạng “Chuyển đổi nhiệt độ giữa thang Celsius và Kelvin”. Phát biểu nào sau đây đúng?
\choice
{Nhiệt lượng và nhiệt độ là hai đại lượng hoàn toàn giống nhau.}
{Mọi quá trình nhiệt học đều không phụ thuộc điều kiện thực hiện.}
{\True $T(K)=t(°C)+273,15.$}
{0 ° $C$ tương ứng 0 K.}
\loigiai{
Phát biểu đúng là: T(K)=t(°C)+273,15. Các phương án còn lại trái với mô hình, định nghĩa hoặc hệ thức nhiệt học tương ứng.
}
\end{ex}

% ===== Câu 71 =====
% ID: L12C1B3-38-TLN
% Mức: VDC

% ===== Câu 72 =====
% ID: L12C1B3-38-TN
% Mức: TH
\begin{ex}
% ID: L12C1B3-38-TN
% Mức: TH
Câu 10 thuộc dạng “Chuyển đổi nhiệt độ giữa thang Celsius và Kelvin”. Phát biểu nào sau đây đúng?
\choice
{Nhiệt lượng và nhiệt độ là hai đại lượng hoàn toàn giống nhau.}
{Mọi quá trình nhiệt học đều không phụ thuộc điều kiện thực hiện.}
{\True $T(K)=t(°C)+273,15.$}
{0 ° $C$ tương ứng 0 K.}
\loigiai{
Phát biểu đúng là: T(K)=t(°C)+273,15. Các phương án còn lại trái với mô hình, định nghĩa hoặc hệ thức nhiệt học tương ứng.
}
\end{ex}

% ===== Câu 73 =====
% ID: L12C1B3-39-TLN
% Mức: VDC

% ===== Câu 74 =====
% ID: L12C1B3-39-TN
% Mức: VD
\begin{ex}
% ID: L12C1B3-39-TN
% Mức: VD
Câu 3 thuộc dạng “Chuyển đổi nhiệt độ giữa thang Celsius và Kelvin”. Phát biểu nào sau đây đúng?
\choice
{Mọi quá trình nhiệt học đều không phụ thuộc điều kiện thực hiện.}
{\True $T(K)=t(°C)+273,15.$}
{0 ° $C$ tương ứng 0 K.}
{Nhiệt lượng và nhiệt độ là hai đại lượng hoàn toàn giống nhau.}
\loigiai{
Phát biểu đúng là: T(K)=t(°C)+273,15. Các phương án còn lại trái với mô hình, định nghĩa hoặc hệ thức nhiệt học tương ứng.
}
\end{ex}

% ===== Câu 75 =====
% ID: L12C1B3-40-TN
% Mức: VD
\begin{ex}
% ID: L12C1B3-40-TN
% Mức: VD
Câu 7 thuộc dạng “Chuyển đổi nhiệt độ giữa thang Celsius và Kelvin”. Phát biểu nào sau đây đúng?
\choice
{Mọi quá trình nhiệt học đều không phụ thuộc điều kiện thực hiện.}
{\True $T(K)=t(°C)+273,15.$}
{0 ° $C$ tương ứng 0 K.}
{Nhiệt lượng và nhiệt độ là hai đại lượng hoàn toàn giống nhau.}
\loigiai{
Phát biểu đúng là: T(K)=t(°C)+273,15. Các phương án còn lại trái với mô hình, định nghĩa hoặc hệ thức nhiệt học tương ứng.
}
\end{ex}

% ===== Câu 76 =====
% ID: L12C1B3-41-TN
% Mức: VDC
\begin{ex}
% ID: L12C1B3-41-TN
% Mức: VDC
Câu 4 thuộc dạng “Chuyển đổi nhiệt độ giữa thang Celsius và Kelvin”. Phát biểu nào sau đây đúng?
\choice
{\True $T(K)=t(°C)+273,15.$}
{0 ° $C$ tương ứng 0 K.}
{Nhiệt lượng và nhiệt độ là hai đại lượng hoàn toàn giống nhau.}
{Mọi quá trình nhiệt học đều không phụ thuộc điều kiện thực hiện.}
\loigiai{
Phát biểu đúng là: T(K)=t(°C)+273,15. Các phương án còn lại trái với mô hình, định nghĩa hoặc hệ thức nhiệt học tương ứng.
}
\end{ex}

% ===== Câu 77 =====
% ID: L12C1B3-42-TN
% Mức: VDC
\begin{ex}
% ID: L12C1B3-42-TN
% Mức: VDC
Câu 8 thuộc dạng “Chuyển đổi nhiệt độ giữa thang Celsius và Kelvin”. Phát biểu nào sau đây đúng?
\choice
{\True $T(K)=t(°C)+273,15.$}
{0 ° $C$ tương ứng 0 K.}
{Nhiệt lượng và nhiệt độ là hai đại lượng hoàn toàn giống nhau.}
{Mọi quá trình nhiệt học đều không phụ thuộc điều kiện thực hiện.}
\loigiai{
Phát biểu đúng là: T(K)=t(°C)+273,15. Các phương án còn lại trái với mô hình, định nghĩa hoặc hệ thức nhiệt học tương ứng.
}
\end{ex}

% ===== Câu 78 =====
% ID: L12C1B3-43-TN
% Mức: NB
\dangbt{Lý thuyết về nhiệt độ và cân bằng nhiệt}
\begin{ex}
% ID: L12C1B3-43-TN
% Mức: NB
Câu 1 thuộc dạng “Lí thuyết về nhiệt độ, cân bằng nhiệt và độ không tuyệt đối”. Phát biểu nào sau đây đúng?
\choice
{Có thể đạt nhiệt độ thấp hơn 0 K.}
{Nhiệt lượng và nhiệt độ là hai đại lượng hoàn toàn giống nhau.}
{Mọi quá trình nhiệt học đều không phụ thuộc điều kiện thực hiện.}
{\True Hai vật cân bằng nhiệt có cùng nhiệt độ.}
\loigiai{
Phát biểu đúng là: Hai vật cân bằng nhiệt có cùng nhiệt độ. Các phương án còn lại trái với mô hình, định nghĩa hoặc hệ thức nhiệt học tương ứng.
}
\end{ex}

% ===== Câu 79 =====
% ID: L12C1B3-44-TN
% Mức: NB
\begin{ex}
% ID: L12C1B3-44-TN
% Mức: NB
Câu 5 thuộc dạng “Lí thuyết về nhiệt độ, cân bằng nhiệt và độ không tuyệt đối”. Phát biểu nào sau đây đúng?
\choice
{Có thể đạt nhiệt độ thấp hơn 0 K.}
{Nhiệt lượng và nhiệt độ là hai đại lượng hoàn toàn giống nhau.}
{Mọi quá trình nhiệt học đều không phụ thuộc điều kiện thực hiện.}
{\True Hai vật cân bằng nhiệt có cùng nhiệt độ.}
\loigiai{
Phát biểu đúng là: Hai vật cân bằng nhiệt có cùng nhiệt độ. Các phương án còn lại trái với mô hình, định nghĩa hoặc hệ thức nhiệt học tương ứng.
}
\end{ex}

% ===== Câu 80 =====
% ID: L12C1B3-45-TN
% Mức: NB
\begin{ex}
% ID: L12C1B3-45-TN
% Mức: NB
Câu 9 thuộc dạng “Lí thuyết về nhiệt độ, cân bằng nhiệt và độ không tuyệt đối”. Phát biểu nào sau đây đúng?
\choice
{Có thể đạt nhiệt độ thấp hơn 0 K.}
{Nhiệt lượng và nhiệt độ là hai đại lượng hoàn toàn giống nhau.}
{Mọi quá trình nhiệt học đều không phụ thuộc điều kiện thực hiện.}
{\True Hai vật cân bằng nhiệt có cùng nhiệt độ.}
\loigiai{
Phát biểu đúng là: Hai vật cân bằng nhiệt có cùng nhiệt độ. Các phương án còn lại trái với mô hình, định nghĩa hoặc hệ thức nhiệt học tương ứng.
}
\end{ex}

% ===== Câu 81 =====
% ID: L12C1B3-46-TN
% Mức: TH
\begin{ex}
% ID: L12C1B3-46-TN
% Mức: TH
Câu 2 thuộc dạng “Lí thuyết về nhiệt độ, cân bằng nhiệt và độ không tuyệt đối”. Phát biểu nào sau đây đúng?
\choice
{Nhiệt lượng và nhiệt độ là hai đại lượng hoàn toàn giống nhau.}
{Mọi quá trình nhiệt học đều không phụ thuộc điều kiện thực hiện.}
{\True Hai vật cân bằng nhiệt có cùng nhiệt độ.}
{Có thể đạt nhiệt độ thấp hơn 0 K.}
\loigiai{
Phát biểu đúng là: Hai vật cân bằng nhiệt có cùng nhiệt độ. Các phương án còn lại trái với mô hình, định nghĩa hoặc hệ thức nhiệt học tương ứng.
}
\end{ex}

% ===== Câu 82 =====
% ID: L12C1B3-47-TN
% Mức: TH
\begin{ex}
% ID: L12C1B3-47-TN
% Mức: TH
Câu 6 thuộc dạng “Lí thuyết về nhiệt độ, cân bằng nhiệt và độ không tuyệt đối”. Phát biểu nào sau đây đúng?
\choice
{Nhiệt lượng và nhiệt độ là hai đại lượng hoàn toàn giống nhau.}
{Mọi quá trình nhiệt học đều không phụ thuộc điều kiện thực hiện.}
{\True Hai vật cân bằng nhiệt có cùng nhiệt độ.}
{Có thể đạt nhiệt độ thấp hơn 0 K.}
\loigiai{
Phát biểu đúng là: Hai vật cân bằng nhiệt có cùng nhiệt độ. Các phương án còn lại trái với mô hình, định nghĩa hoặc hệ thức nhiệt học tương ứng.
}
\end{ex}

% ===== Câu 83 =====
% ID: L12C1B3-48-TN
% Mức: TH
\begin{ex}
% ID: L12C1B3-48-TN
% Mức: TH
Câu 10 thuộc dạng “Lí thuyết về nhiệt độ, cân bằng nhiệt và độ không tuyệt đối”. Phát biểu nào sau đây đúng?
\choice
{Nhiệt lượng và nhiệt độ là hai đại lượng hoàn toàn giống nhau.}
{Mọi quá trình nhiệt học đều không phụ thuộc điều kiện thực hiện.}
{\True Hai vật cân bằng nhiệt có cùng nhiệt độ.}
{Có thể đạt nhiệt độ thấp hơn 0 K.}
\loigiai{
Phát biểu đúng là: Hai vật cân bằng nhiệt có cùng nhiệt độ. Các phương án còn lại trái với mô hình, định nghĩa hoặc hệ thức nhiệt học tương ứng.
}
\end{ex}

% ===== Câu 84 =====
% ID: L12C1B3-49-TN
% Mức: VD
\begin{ex}
% ID: L12C1B3-49-TN
% Mức: VD
Câu 3 thuộc dạng “Lí thuyết về nhiệt độ, cân bằng nhiệt và độ không tuyệt đối”. Phát biểu nào sau đây đúng?
\choice
{Mọi quá trình nhiệt học đều không phụ thuộc điều kiện thực hiện.}
{\True Hai vật cân bằng nhiệt có cùng nhiệt độ.}
{Có thể đạt nhiệt độ thấp hơn 0 K.}
{Nhiệt lượng và nhiệt độ là hai đại lượng hoàn toàn giống nhau.}
\loigiai{
Phát biểu đúng là: Hai vật cân bằng nhiệt có cùng nhiệt độ. Các phương án còn lại trái với mô hình, định nghĩa hoặc hệ thức nhiệt học tương ứng.
}
\end{ex}

% ===== Câu 85 =====
% ID: L12C1B3-50-TN
% Mức: VD
\begin{ex}
% ID: L12C1B3-50-TN
% Mức: VD
Câu 7 thuộc dạng “Lí thuyết về nhiệt độ, cân bằng nhiệt và độ không tuyệt đối”. Phát biểu nào sau đây đúng?
\choice
{Mọi quá trình nhiệt học đều không phụ thuộc điều kiện thực hiện.}
{\True Hai vật cân bằng nhiệt có cùng nhiệt độ.}
{Có thể đạt nhiệt độ thấp hơn 0 K.}
{Nhiệt lượng và nhiệt độ là hai đại lượng hoàn toàn giống nhau.}
\loigiai{
Phát biểu đúng là: Hai vật cân bằng nhiệt có cùng nhiệt độ. Các phương án còn lại trái với mô hình, định nghĩa hoặc hệ thức nhiệt học tương ứng.
}
\end{ex}

% ===== Câu 86 =====
% ID: L12C1B3-51-TN
% Mức: VDC
\begin{ex}
% ID: L12C1B3-51-TN
% Mức: VDC
Câu 4 thuộc dạng “Lí thuyết về nhiệt độ, cân bằng nhiệt và độ không tuyệt đối”. Phát biểu nào sau đây đúng?
\choice
{\True Hai vật cân bằng nhiệt có cùng nhiệt độ.}
{Có thể đạt nhiệt độ thấp hơn 0 K.}
{Nhiệt lượng và nhiệt độ là hai đại lượng hoàn toàn giống nhau.}
{Mọi quá trình nhiệt học đều không phụ thuộc điều kiện thực hiện.}
\loigiai{
Phát biểu đúng là: Hai vật cân bằng nhiệt có cùng nhiệt độ. Các phương án còn lại trái với mô hình, định nghĩa hoặc hệ thức nhiệt học tương ứng.
}
\end{ex}

% ===== Câu 87 =====
% ID: L12C1B3-52-TN
% Mức: VDC
\begin{ex}
% ID: L12C1B3-52-TN
% Mức: VDC
Câu 8 thuộc dạng “Lí thuyết về nhiệt độ, cân bằng nhiệt và độ không tuyệt đối”. Phát biểu nào sau đây đúng?
\choice
{\True Hai vật cân bằng nhiệt có cùng nhiệt độ.}
{Có thể đạt nhiệt độ thấp hơn 0 K.}
{Nhiệt lượng và nhiệt độ là hai đại lượng hoàn toàn giống nhau.}
{Mọi quá trình nhiệt học đều không phụ thuộc điều kiện thực hiện.}
\loigiai{
Phát biểu đúng là: Hai vật cân bằng nhiệt có cùng nhiệt độ. Các phương án còn lại trái với mô hình, định nghĩa hoặc hệ thức nhiệt học tương ứng.
}
\end{ex}

% ===== Câu 88 =====
% ID: L12C1B3-53-TN
% Mức: NB
\dangbt{Nguyên lý hoạt động và sử dụng nhiệt kế}
\begin{ex}
% ID: L12C1B3-53-TN
% Mức: NB
Câu 1 thuộc dạng “Nguyên lí hoạt động và sử dụng nhiệt kế”. Phát biểu nào sau đây đúng?
\choice
{Có thể dùng mọi nhiệt kế để đo mọi khoảng nhiệt độ.}
{Nhiệt lượng và nhiệt độ là hai đại lượng hoàn toàn giống nhau.}
{Mọi quá trình nhiệt học đều không phụ thuộc điều kiện thực hiện.}
{\True Nhiệt kế hoạt động dựa trên đại lượng vật lí biến đổi xác định theo nhiệt độ.}
\loigiai{
Phát biểu đúng là: Nhiệt kế hoạt động dựa trên đại lượng vật lí biến đổi xác định theo nhiệt độ. Các phương án còn lại trái với mô hình, định nghĩa hoặc hệ thức nhiệt học tương ứng.
}
\end{ex}

% ===== Câu 89 =====
% ID: L12C1B3-54-TN
% Mức: NB
\begin{ex}
% ID: L12C1B3-54-TN
% Mức: NB
Câu 5 thuộc dạng “Nguyên lí hoạt động và sử dụng nhiệt kế”. Phát biểu nào sau đây đúng?
\choice
{Có thể dùng mọi nhiệt kế để đo mọi khoảng nhiệt độ.}
{Nhiệt lượng và nhiệt độ là hai đại lượng hoàn toàn giống nhau.}
{Mọi quá trình nhiệt học đều không phụ thuộc điều kiện thực hiện.}
{\True Nhiệt kế hoạt động dựa trên đại lượng vật lí biến đổi xác định theo nhiệt độ.}
\loigiai{
Phát biểu đúng là: Nhiệt kế hoạt động dựa trên đại lượng vật lí biến đổi xác định theo nhiệt độ. Các phương án còn lại trái với mô hình, định nghĩa hoặc hệ thức nhiệt học tương ứng.
}
\end{ex}

% ===== Câu 90 =====
% ID: L12C1B3-55-TN
% Mức: NB
\begin{ex}
% ID: L12C1B3-55-TN
% Mức: NB
Câu 9 thuộc dạng “Nguyên lí hoạt động và sử dụng nhiệt kế”. Phát biểu nào sau đây đúng?
\choice
{Có thể dùng mọi nhiệt kế để đo mọi khoảng nhiệt độ.}
{Nhiệt lượng và nhiệt độ là hai đại lượng hoàn toàn giống nhau.}
{Mọi quá trình nhiệt học đều không phụ thuộc điều kiện thực hiện.}
{\True Nhiệt kế hoạt động dựa trên đại lượng vật lí biến đổi xác định theo nhiệt độ.}
\loigiai{
Phát biểu đúng là: Nhiệt kế hoạt động dựa trên đại lượng vật lí biến đổi xác định theo nhiệt độ. Các phương án còn lại trái với mô hình, định nghĩa hoặc hệ thức nhiệt học tương ứng.
}
\end{ex}

% ===== Câu 91 =====
% ID: L12C1B3-56-TN
% Mức: TH
\begin{ex}
% ID: L12C1B3-56-TN
% Mức: TH
Câu 2 thuộc dạng “Nguyên lí hoạt động và sử dụng nhiệt kế”. Phát biểu nào sau đây đúng?
\choice
{Nhiệt lượng và nhiệt độ là hai đại lượng hoàn toàn giống nhau.}
{Mọi quá trình nhiệt học đều không phụ thuộc điều kiện thực hiện.}
{\True Nhiệt kế hoạt động dựa trên đại lượng vật lí biến đổi xác định theo nhiệt độ.}
{Có thể dùng mọi nhiệt kế để đo mọi khoảng nhiệt độ.}
\loigiai{
Phát biểu đúng là: Nhiệt kế hoạt động dựa trên đại lượng vật lí biến đổi xác định theo nhiệt độ. Các phương án còn lại trái với mô hình, định nghĩa hoặc hệ thức nhiệt học tương ứng.
}
\end{ex}

% ===== Câu 92 =====
% ID: L12C1B3-57-TN
% Mức: TH
\begin{ex}
% ID: L12C1B3-57-TN
% Mức: TH
Câu 6 thuộc dạng “Nguyên lí hoạt động và sử dụng nhiệt kế”. Phát biểu nào sau đây đúng?
\choice
{Nhiệt lượng và nhiệt độ là hai đại lượng hoàn toàn giống nhau.}
{Mọi quá trình nhiệt học đều không phụ thuộc điều kiện thực hiện.}
{\True Nhiệt kế hoạt động dựa trên đại lượng vật lí biến đổi xác định theo nhiệt độ.}
{Có thể dùng mọi nhiệt kế để đo mọi khoảng nhiệt độ.}
\loigiai{
Phát biểu đúng là: Nhiệt kế hoạt động dựa trên đại lượng vật lí biến đổi xác định theo nhiệt độ. Các phương án còn lại trái với mô hình, định nghĩa hoặc hệ thức nhiệt học tương ứng.
}
\end{ex}

% ===== Câu 93 =====
% ID: L12C1B3-58-TN
% Mức: TH
\begin{ex}
% ID: L12C1B3-58-TN
% Mức: TH
Câu 10 thuộc dạng “Nguyên lí hoạt động và sử dụng nhiệt kế”. Phát biểu nào sau đây đúng?
\choice
{Nhiệt lượng và nhiệt độ là hai đại lượng hoàn toàn giống nhau.}
{Mọi quá trình nhiệt học đều không phụ thuộc điều kiện thực hiện.}
{\True Nhiệt kế hoạt động dựa trên đại lượng vật lí biến đổi xác định theo nhiệt độ.}
{Có thể dùng mọi nhiệt kế để đo mọi khoảng nhiệt độ.}
\loigiai{
Phát biểu đúng là: Nhiệt kế hoạt động dựa trên đại lượng vật lí biến đổi xác định theo nhiệt độ. Các phương án còn lại trái với mô hình, định nghĩa hoặc hệ thức nhiệt học tương ứng.
}
\end{ex}

% ===== Câu 94 =====
% ID: L12C1B3-59-TN
% Mức: VD
\begin{ex}
% ID: L12C1B3-59-TN
% Mức: VD
Câu 3 thuộc dạng “Nguyên lí hoạt động và sử dụng nhiệt kế”. Phát biểu nào sau đây đúng?
\choice
{Mọi quá trình nhiệt học đều không phụ thuộc điều kiện thực hiện.}
{\True Nhiệt kế hoạt động dựa trên đại lượng vật lí biến đổi xác định theo nhiệt độ.}
{Có thể dùng mọi nhiệt kế để đo mọi khoảng nhiệt độ.}
{Nhiệt lượng và nhiệt độ là hai đại lượng hoàn toàn giống nhau.}
\loigiai{
Phát biểu đúng là: Nhiệt kế hoạt động dựa trên đại lượng vật lí biến đổi xác định theo nhiệt độ. Các phương án còn lại trái với mô hình, định nghĩa hoặc hệ thức nhiệt học tương ứng.
}
\end{ex}

% ===== Câu 95 =====
% ID: L12C1B3-60-TN
% Mức: VD
\begin{ex}
% ID: L12C1B3-60-TN
% Mức: VD
Câu 7 thuộc dạng “Nguyên lí hoạt động và sử dụng nhiệt kế”. Phát biểu nào sau đây đúng?
\choice
{Mọi quá trình nhiệt học đều không phụ thuộc điều kiện thực hiện.}
{\True Nhiệt kế hoạt động dựa trên đại lượng vật lí biến đổi xác định theo nhiệt độ.}
{Có thể dùng mọi nhiệt kế để đo mọi khoảng nhiệt độ.}
{Nhiệt lượng và nhiệt độ là hai đại lượng hoàn toàn giống nhau.}
\loigiai{
Phát biểu đúng là: Nhiệt kế hoạt động dựa trên đại lượng vật lí biến đổi xác định theo nhiệt độ. Các phương án còn lại trái với mô hình, định nghĩa hoặc hệ thức nhiệt học tương ứng.
}
\end{ex}

% ===== Câu 96 =====
% ID: L12C1B3-61-TN
% Mức: VDC
\begin{ex}
% ID: L12C1B3-61-TN
% Mức: VDC
Câu 4 thuộc dạng “Nguyên lí hoạt động và sử dụng nhiệt kế”. Phát biểu nào sau đây đúng?
\choice
{\True Nhiệt kế hoạt động dựa trên đại lượng vật lí biến đổi xác định theo nhiệt độ.}
{Có thể dùng mọi nhiệt kế để đo mọi khoảng nhiệt độ.}
{Nhiệt lượng và nhiệt độ là hai đại lượng hoàn toàn giống nhau.}
{Mọi quá trình nhiệt học đều không phụ thuộc điều kiện thực hiện.}
\loigiai{
Phát biểu đúng là: Nhiệt kế hoạt động dựa trên đại lượng vật lí biến đổi xác định theo nhiệt độ. Các phương án còn lại trái với mô hình, định nghĩa hoặc hệ thức nhiệt học tương ứng.
}
\end{ex}

% ===== Câu 97 =====
% ID: L12C1B3-62-TN
% Mức: VDC
\begin{ex}
% ID: L12C1B3-62-TN
% Mức: VDC
Câu 8 thuộc dạng “Nguyên lí hoạt động và sử dụng nhiệt kế”. Phát biểu nào sau đây đúng?
\choice
{\True Nhiệt kế hoạt động dựa trên đại lượng vật lí biến đổi xác định theo nhiệt độ.}
{Có thể dùng mọi nhiệt kế để đo mọi khoảng nhiệt độ.}
{Nhiệt lượng và nhiệt độ là hai đại lượng hoàn toàn giống nhau.}
{Mọi quá trình nhiệt học đều không phụ thuộc điều kiện thực hiện.}
\loigiai{
Phát biểu đúng là: Nhiệt kế hoạt động dựa trên đại lượng vật lí biến đổi xác định theo nhiệt độ. Các phương án còn lại trái với mô hình, định nghĩa hoặc hệ thức nhiệt học tương ứng.
}
\end{ex}

% ===== Câu 98 =====
% ID: L12C1B3-63-TN
% Mức: NB
\dangbt{Thiết lập và chuyển đổi thang nhiệt độ tự chế}
\begin{ex}
% ID: L12C1B3-63-TN
% Mức: NB
Câu 1 thuộc dạng “Thiết lập và chuyển đổi thang nhiệt độ mới”. Phát biểu nào sau đây đúng?
\choice
{Khoảng chia của mọi thang nhiệt độ đều bằng nhau.}
{Nhiệt lượng và nhiệt độ là hai đại lượng hoàn toàn giống nhau.}
{Mọi quá trình nhiệt học đều không phụ thuộc điều kiện thực hiện.}
{\True Thang tuyến tính được xác định từ hai mốc và phép nội suy.}
\loigiai{
Phát biểu đúng là: Thang tuyến tính được xác định từ hai mốc và phép nội suy. Các phương án còn lại trái với mô hình, định nghĩa hoặc hệ thức nhiệt học tương ứng.
}
\end{ex}

% ===== Câu 99 =====
% ID: L12C1B3-64-TN
% Mức: NB
\begin{ex}
% ID: L12C1B3-64-TN
% Mức: NB
Câu 5 thuộc dạng “Thiết lập và chuyển đổi thang nhiệt độ mới”. Phát biểu nào sau đây đúng?
\choice
{Khoảng chia của mọi thang nhiệt độ đều bằng nhau.}
{Nhiệt lượng và nhiệt độ là hai đại lượng hoàn toàn giống nhau.}
{Mọi quá trình nhiệt học đều không phụ thuộc điều kiện thực hiện.}
{\True Thang tuyến tính được xác định từ hai mốc và phép nội suy.}
\loigiai{
Phát biểu đúng là: Thang tuyến tính được xác định từ hai mốc và phép nội suy. Các phương án còn lại trái với mô hình, định nghĩa hoặc hệ thức nhiệt học tương ứng.
}
\end{ex}

% ===== Câu 100 =====
% ID: L12C1B3-65-TN
% Mức: NB
\begin{ex}
% ID: L12C1B3-65-TN
% Mức: NB
Câu 9 thuộc dạng “Thiết lập và chuyển đổi thang nhiệt độ mới”. Phát biểu nào sau đây đúng?
\choice
{Khoảng chia của mọi thang nhiệt độ đều bằng nhau.}
{Nhiệt lượng và nhiệt độ là hai đại lượng hoàn toàn giống nhau.}
{Mọi quá trình nhiệt học đều không phụ thuộc điều kiện thực hiện.}
{\True Thang tuyến tính được xác định từ hai mốc và phép nội suy.}
\loigiai{
Phát biểu đúng là: Thang tuyến tính được xác định từ hai mốc và phép nội suy. Các phương án còn lại trái với mô hình, định nghĩa hoặc hệ thức nhiệt học tương ứng.
}
\end{ex}

% ===== Câu 101 =====
% ID: L12C1B3-66-TN
% Mức: TH
\begin{ex}
% ID: L12C1B3-66-TN
% Mức: TH
Câu 2 thuộc dạng “Thiết lập và chuyển đổi thang nhiệt độ mới”. Phát biểu nào sau đây đúng?
\choice
{Nhiệt lượng và nhiệt độ là hai đại lượng hoàn toàn giống nhau.}
{Mọi quá trình nhiệt học đều không phụ thuộc điều kiện thực hiện.}
{\True Thang tuyến tính được xác định từ hai mốc và phép nội suy.}
{Khoảng chia của mọi thang nhiệt độ đều bằng nhau.}
\loigiai{
Phát biểu đúng là: Thang tuyến tính được xác định từ hai mốc và phép nội suy. Các phương án còn lại trái với mô hình, định nghĩa hoặc hệ thức nhiệt học tương ứng.
}
\end{ex}

% ===== Câu 102 =====
% ID: L12C1B3-67-TN
% Mức: TH
\begin{ex}
% ID: L12C1B3-67-TN
% Mức: TH
Câu 6 thuộc dạng “Thiết lập và chuyển đổi thang nhiệt độ mới”. Phát biểu nào sau đây đúng?
\choice
{Nhiệt lượng và nhiệt độ là hai đại lượng hoàn toàn giống nhau.}
{Mọi quá trình nhiệt học đều không phụ thuộc điều kiện thực hiện.}
{\True Thang tuyến tính được xác định từ hai mốc và phép nội suy.}
{Khoảng chia của mọi thang nhiệt độ đều bằng nhau.}
\loigiai{
Phát biểu đúng là: Thang tuyến tính được xác định từ hai mốc và phép nội suy. Các phương án còn lại trái với mô hình, định nghĩa hoặc hệ thức nhiệt học tương ứng.
}
\end{ex}

% ===== Câu 103 =====
% ID: L12C1B3-68-TN
% Mức: TH
\begin{ex}
% ID: L12C1B3-68-TN
% Mức: TH
Câu 10 thuộc dạng “Thiết lập và chuyển đổi thang nhiệt độ mới”. Phát biểu nào sau đây đúng?
\choice
{Nhiệt lượng và nhiệt độ là hai đại lượng hoàn toàn giống nhau.}
{Mọi quá trình nhiệt học đều không phụ thuộc điều kiện thực hiện.}
{\True Thang tuyến tính được xác định từ hai mốc và phép nội suy.}
{Khoảng chia của mọi thang nhiệt độ đều bằng nhau.}
\loigiai{
Phát biểu đúng là: Thang tuyến tính được xác định từ hai mốc và phép nội suy. Các phương án còn lại trái với mô hình, định nghĩa hoặc hệ thức nhiệt học tương ứng.
}
\end{ex}

% ===== Câu 104 =====
% ID: L12C1B3-69-TN
% Mức: VD
\begin{ex}
% ID: L12C1B3-69-TN
% Mức: VD
Câu 3 thuộc dạng “Thiết lập và chuyển đổi thang nhiệt độ mới”. Phát biểu nào sau đây đúng?
\choice
{Mọi quá trình nhiệt học đều không phụ thuộc điều kiện thực hiện.}
{\True Thang tuyến tính được xác định từ hai mốc và phép nội suy.}
{Khoảng chia của mọi thang nhiệt độ đều bằng nhau.}
{Nhiệt lượng và nhiệt độ là hai đại lượng hoàn toàn giống nhau.}
\loigiai{
Phát biểu đúng là: Thang tuyến tính được xác định từ hai mốc và phép nội suy. Các phương án còn lại trái với mô hình, định nghĩa hoặc hệ thức nhiệt học tương ứng.
}
\end{ex}

% ===== Câu 105 =====
% ID: L12C1B3-70-TN
% Mức: VD
\begin{ex}
% ID: L12C1B3-70-TN
% Mức: VD
Câu 7 thuộc dạng “Thiết lập và chuyển đổi thang nhiệt độ mới”. Phát biểu nào sau đây đúng?
\choice
{Mọi quá trình nhiệt học đều không phụ thuộc điều kiện thực hiện.}
{\True Thang tuyến tính được xác định từ hai mốc và phép nội suy.}
{Khoảng chia của mọi thang nhiệt độ đều bằng nhau.}
{Nhiệt lượng và nhiệt độ là hai đại lượng hoàn toàn giống nhau.}
\loigiai{
Phát biểu đúng là: Thang tuyến tính được xác định từ hai mốc và phép nội suy. Các phương án còn lại trái với mô hình, định nghĩa hoặc hệ thức nhiệt học tương ứng.
}
\end{ex}

% ===== Câu 106 =====
% ID: L12C1B3-71-TN
% Mức: VDC
\begin{ex}
% ID: L12C1B3-71-TN
% Mức: VDC
Câu 4 thuộc dạng “Thiết lập và chuyển đổi thang nhiệt độ mới”. Phát biểu nào sau đây đúng?
\choice
{\True Thang tuyến tính được xác định từ hai mốc và phép nội suy.}
{Khoảng chia của mọi thang nhiệt độ đều bằng nhau.}
{Nhiệt lượng và nhiệt độ là hai đại lượng hoàn toàn giống nhau.}
{Mọi quá trình nhiệt học đều không phụ thuộc điều kiện thực hiện.}
\loigiai{
Phát biểu đúng là: Thang tuyến tính được xác định từ hai mốc và phép nội suy. Các phương án còn lại trái với mô hình, định nghĩa hoặc hệ thức nhiệt học tương ứng.
}
\end{ex}

% ===== Câu 107 =====
% ID: L12C1B3-72-TN
% Mức: VDC
\begin{ex}
% ID: L12C1B3-72-TN
% Mức: VDC
Câu 8 thuộc dạng “Thiết lập và chuyển đổi thang nhiệt độ mới”. Phát biểu nào sau đây đúng?
\choice
{\True Thang tuyến tính được xác định từ hai mốc và phép nội suy.}
{Khoảng chia của mọi thang nhiệt độ đều bằng nhau.}
{Nhiệt lượng và nhiệt độ là hai đại lượng hoàn toàn giống nhau.}
{Mọi quá trình nhiệt học đều không phụ thuộc điều kiện thực hiện.}
\loigiai{
Phát biểu đúng là: Thang tuyến tính được xác định từ hai mốc và phép nội suy. Các phương án còn lại trái với mô hình, định nghĩa hoặc hệ thức nhiệt học tương ứng.
}
\end{ex}
\dangbt{Lý thuyết về nhiệt độ và cân bằng nhiệt}
\begin{ex}
Nhiệt độ cho biết điều gì về hai vật khi chúng tiếp xúc nhau?
\choice
{Khối lượng của hai vật}
{\True Trạng thái cân bằng nhiệt}
{Hình dạng của hai vật}
{Màu sắc của hai vật}
\loigiai{Nhiệt độ cho biết trạng thái cân bằng nhiệt giữa hai vật khi chúng tiếp xúc nhau.}
\end{ex}

\dangbt{Lý thuyết về nhiệt độ và cân bằng nhiệt}
\begin{ex}
Khi hai vật có nhiệt độ bằng nhau, điều gì sẽ xảy ra?
\choice
{Vật lớn sẽ truyền nhiệt cho vật nhỏ}
{\True Không có sự truyền nhiệt}
{Vật nhỏ sẽ truyền nhiệt cho vật lớn}
{Cả hai vật đều mất nhiệt}
\loigiai{Khi hai vật có nhiệt độ bằng nhau, không có sự truyền nhiệt giữa chúng.}
\end{ex}

\dangbt{Khái niệm nhiệt độ}
\begin{ex}
Nhiệt độ thấp nhất mà các vật có thể đạt được được gọi là gì?
\choice
{Nhiệt độ nóng chảy}
{Nhiệt độ điểm ba}
{\True Độ không tuyệt đối}
{Nhiệt độ sôi}
\loigiai{Nhiệt độ thấp nhất mà các vật có thể đạt được được gọi là độ không tuyệt đối.}
\end{ex}

\dangbt{Khái niệm nhiệt độ}
\begin{ex}
Một vật có nhiệt độ $-273.15^\circ C$ được gọi là gì?
\choice
{Nhiệt độ nóng chảy}
{Nhiệt độ điểm ba}
{\True Độ không tuyệt đối}
{Nhiệt độ sôi}
\loigiai{Nhiệt độ $-273.15^\circ C$ là độ không tuyệt đối, nhiệt độ thấp nhất mà một vật có thể đạt được.}
\end{ex}

\dangbt{Khái niệm nhiệt độ}
\begin{ex}
Nhiệt độ thấp nhất mà một vật có thể đạt được là?
\choice
{\True $-273.15^\circ C$}
{$0^\circ C$}
{$-100^\circ C$}
{$100K$}
\loigiai{Nhiệt độ thấp nhất mà một vật có thể đạt được là $-273.15^\circ C$, còn gọi là độ không tuyệt đối.}
\end{ex}

\dangbt{Thang nhiệt độ và nhiệt kế}
\begin{ex}
Nhiệt độ nóng chảy của nước đá tinh khiết là bao nhiêu độ Celsius?
\choice
{$100^\circ C$}
{$50^\circ C$}
{\True $0^\circ C$}
{$-10^\circ C$}
\loigiai{Nhiệt độ nóng chảy của nước đá tinh khiết là $0^\circ C$.}
\end{ex}

\dangbt{Thang nhiệt độ và nhiệt kế}
\begin{ex}
Thang nhiệt độ nào sử dụng độ không tuyệt đối làm mốc?
\choice
{Celsius}
{Fahrenheit}
{\True Kelvin}
{Rankine}
\loigiai{Thang nhiệt độ Kelvin sử dụng độ không tuyệt đối ($0K$) làm mốc.}
\end{ex}

\dangbt{Thang nhiệt độ và nhiệt kế}
\begin{ex}
Nhiệt độ sôi của nước tinh khiết ở áp suất $1$ atm là bao nhiêu độ Fahrenheit?
\choice
{$32^\circ F$}
{$100^\circ F$}
{\True $212^\circ F$}
{$273^\circ F$}
\loigiai{Nhiệt độ sôi của nước tinh khiết ở áp suất $1$ atm là $212^\circ F$.}
\end{ex}

\dangbt{Thang nhiệt độ và nhiệt kế}
\begin{ex}
Đơn vị của nhiệt độ trong thang đo Kelvin là gì?
\choice
{$^\circ C$}
{$^\circ F$}
{\True K}
{J}
\loigiai{Đơn vị của nhiệt độ trong thang đo Kelvin là K (Kelvin).}
\end{ex}

\dangbt{Thang nhiệt độ và nhiệt kế}
\begin{ex}
Điểm ba của nước có nhiệt độ bao nhiêu độ Kelvin?
\choice
{$100K$}
{\True $273.15K$}
{$373.15K$}
{$273K$}
\loigiai{Điểm ba của nước có nhiệt độ $273.15K$.}
\end{ex}

\dangbt{Thang nhiệt độ và nhiệt kế}
\begin{ex}
Thang nhiệt độ nào được sử dụng phổ biến ở các nước phương Tây?
\choice
{Celsius}
{Kelvin}
{\True Fahrenheit}
{Rankine}
\loigiai{Thang nhiệt độ Fahrenheit được sử dụng phổ biến ở các nước phương Tây.}
\end{ex}

\dangbt{Thang nhiệt độ và nhiệt kế}
\begin{ex}
Nhiệt độ điểm ba của nước trong thang Celsius là bao nhiêu?
\choice
{$0^\circ C$}
{$100^\circ C$}
{\True $0.01^\circ C$}
{$273.15^\circ C$}
\loigiai{Nhiệt độ điểm ba của nước trong thang Celsius là $0.01^\circ C$.}
\end{ex}

\dangbt{Thang nhiệt độ và nhiệt kế}
\begin{ex}
Thang đo nhiệt độ nào chia khoảng giữa nhiệt độ nóng chảy và nhiệt độ sôi của nước thành $100$ khoảng bằng nhau?
\choice
{Kelvin}
{Fahrenheit}
{\True Celsius}
{Rankine}
\loigiai{Thang đo Celsius chia khoảng giữa nhiệt độ nóng chảy và nhiệt độ sôi của nước thành $100$ khoảng bằng nhau.}
\end{ex}

\dangbt{Thang nhiệt độ và nhiệt kế}
\begin{ex}
Khi đo nhiệt độ của nước đang sôi bằng nhiệt kế Celsius, kết quả sẽ là bao nhiêu?
\choice
{$0^\circ C$}
{$50^\circ C$}
{\True $100^\circ C$}
{$150^\circ C$}
\loigiai{Nhiệt độ sôi của nước tinh khiết ở áp suất $1$ atm là $100^\circ C$.}
\end{ex}

\dangbt{Nguyên lý hoạt động và sử dụng nhiệt kế}
\begin{ex}
Nhiệt kế nào sử dụng sự nở dài của cột chất lỏng trong ống thủy tinh để đo nhiệt độ?
\choice
{Nhiệt kế kim loại}
{\True Nhiệt kế rượu}
{Nhiệt kế hồng ngoại}
{Nhiệt kế khí}
\loigiai{Nhiệt kế rượu sử dụng sự nở dài của cột rượu trong ống thủy tinh để đo nhiệt độ.}
\end{ex}

\dangbt{Nguyên lý hoạt động và sử dụng nhiệt kế}
\begin{ex}
Cảm biến hồng ngoại đo gì để tính toán nhiệt độ?
\choice
{Khối lượng của vật}
{\True Mức năng lượng bức xạ}
{Tốc độ di chuyển của vật}
{Kích thước của vật}
\loigiai{Cảm biến hồng ngoại đo mức năng lượng bức xạ để tính toán nhiệt độ.}
\end{ex}

\dangbt{Nguyên lý hoạt động và sử dụng nhiệt kế}
\begin{ex}
Nhiệt kế kim loại đo nhiệt độ dựa trên nguyên lý gì?
\choice
{Sự nở dài của cột chất lỏng}
{Sự thay đổi điện trở của kim loại}
{\True Sự nở dài của thanh kim loại mỏng}
{Sự thay đổi áp suất của khí}
\loigiai{Nhiệt kế kim loại đo nhiệt độ dựa trên sự nở dài của thanh kim loại mỏng.}
\end{ex}

\dangbt{Nguyên lý hoạt động và sử dụng nhiệt kế}
\begin{ex}
Khi nhiệt độ tăng, cột chất lỏng trong nhiệt kế sẽ làm gì?
\choice
{Co lại}
{\True Dài ra}
{Không thay đổi}
{Chuyển thành thể rắn}
\loigiai{Khi nhiệt độ tăng, cột chất lỏng trong nhiệt kế sẽ dài ra.}
\end{ex}

\dangbt{Nguyên lý hoạt động và sử dụng nhiệt kế}
\begin{ex}
Trong quá trình đo nhiệt độ bằng nhiệt kế khí, yếu tố nào không đổi?
\choice
{Áp suất của khí}
{Thể tích của khí}
{Khối lượng của khí}
{\True Nhiệt độ của khí}
\loigiai{Trong quá trình đo nhiệt độ bằng nhiệt kế khí, áp suất của khí không đổi.}
\end{ex}

\dangbt{Nguyên lý hoạt động và sử dụng nhiệt kế}
\begin{ex}
Nhiệt kế nào hoạt động dựa trên sự nở vì nhiệt của thể tích khí?
\choice
{Nhiệt kế rượu}
{Nhiệt kế kim loại}
{\True Nhiệt kế khí}
{Nhiệt kế hồng ngoại}
\loigiai{Nhiệt kế khí hoạt động dựa trên sự nở vì nhiệt của thể tích khí.}
\end{ex}

\dangbt{Nguyên lý hoạt động và sử dụng nhiệt kế}
\begin{ex}
Một nhiệt kế hồng ngoại điện tử đo nhiệt độ dựa trên nguyên lý nào?
\choice
{Sự nở dài của kim loại}
{Sự thay đổi áp suất của khí}
{\True Bức xạ điện từ của vật}
{Sự thay đổi điện trở của kim loại}
\loigiai{Nhiệt kế hồng ngoại điện tử đo nhiệt độ dựa trên bức xạ điện từ phát ra từ vật có nhiệt độ trên $-273^\circ C$.}
\end{ex}

\dangbt{Nguyên lý hoạt động và sử dụng nhiệt kế}
\begin{ex}
Một nhiệt kế kim loại dựa trên sự nở dài của thanh kim loại sẽ cho kết quả chính xác nhất trong điều kiện nào?
\choice
{\True Khi đo nhiệt độ không đổi}
{Khi đo nhiệt độ thay đổi nhanh chóng}
{Khi đo nhiệt độ của chất lỏng}
{Khi đo nhiệt độ của khí}
\loigiai{Nhiệt kế kim loại dựa trên sự nở dài của thanh kim loại sẽ cho kết quả chính xác nhất khi đo nhiệt độ không đổi. Sự ổn định nhiệt độ giúp thanh kim loại nở dài đến mức cân bằng, phản ánh chính xác nhiệt độ đo được.}
\end{ex}

\dangbt{Chuyển đổi nhiệt độ giữa các thang đo thông dụng}
\begin{ex}
Công thức chuyển đổi từ thang đo Celsius sang Kelvin là gì?
\choice
{\True $T(K)= t(^\circ C)+ 273.15$}
{$T(K)= t(^\circ C)+ 32$}
{$T(K)= t(^\circ C)/ 1.8$}
{$T(K)= t(^\circ C)* 1.8 + 32$}
\loigiai{Công thức chuyển đổi từ thang đo Celsius sang Kelvin là $T(K)= t(^\circ C)+ 273.15$.}
\end{ex}

\dangbt{Chuyển đổi nhiệt độ giữa các thang đo thông dụng}
\begin{ex}
Một vật có nhiệt độ $-40^\circ C$ sẽ tương đương với bao nhiêu độ Fahrenheit?
\choice
{\True $-40^\circ F$}
{$0^\circ F$}
{$32^\circ F$}
{$-72^\circ F$}
\loigiai{Nhiệt độ $-40^\circ C$ tương đương với $-40^\circ F$.}
\end{ex}

\dangbt{Chuyển đổi nhiệt độ giữa các thang đo thông dụng}
\begin{ex}
Nhiệt độ $32^\circ F$ tương đương với bao nhiêu độ Celsius?
\choice
{\True $0^\circ C$}
{$32^\circ C$}
{$100^\circ C$}
{$273.15^\circ C$}
\loigiai{Nhiệt độ $32^\circ F$ tương đương với $0^\circ C$.}
\end{ex}

\dangbt{Chuyển đổi nhiệt độ giữa các thang đo thông dụng}
\begin{ex}
Nếu một vật có nhiệt độ $25^\circ C$, nhiệt độ này tương đương bao nhiêu độ Kelvin?
\choice
{\True $298.15K$}
{$273.15K$}
{$250K$}
{$300K$}
\loigiai{$T(K)= t(^\circ C)+ 273.15$, do đó, $25^\circ C = 298.15K$.}
\end{ex}

\dangbt{Chuyển đổi nhiệt độ giữa các thang đo thông dụng}
\begin{ex}
Nếu nhiệt độ của một vật là $77^\circ F$, nhiệt độ này tương đương bao nhiêu độ Celsius?
\choice
{\True $25^\circ C$}
{$0^\circ C$}
{$20^\circ C$}
{$30^\circ C$}
\loigiai{$T(^\circ C)= (T(^\circ F)- 32)/ 1.8$, do đó, $77^\circ F = 25^\circ C$.}
\end{ex}

\dangbt{Chuyển đổi nhiệt độ giữa các thang đo thông dụng}
\begin{ex}
Để chuyển đổi từ nhiệt độ Celsius sang Fahrenheit, công thức nào được sử dụng?
\choice
{$T(^\circ F)= t(^\circ C)+ 273.15$}
{$T(^\circ F)= t(^\circ C)/ 1.8$}
{\True $T(^\circ F)= t(^\circ C)* 1.8 + 32$}
{$T(^\circ F)= t(^\circ C)- 32 / 1.8$}
\loigiai{Công thức chuyển đổi từ nhiệt độ Celsius sang Fahrenheit là $T(^\circ F)= t(^\circ C)* 1.8 + 32$.}
\end{ex}

\dangbt{Chuyển đổi nhiệt độ giữa các thang đo thông dụng}
\begin{ex}
Nhiệt độ của một vật là $300K$, nhiệt độ này tương đương bao nhiêu độ Celsius?
\choice
{\True $27^\circ C$}
{$30^\circ C$}
{$25^\circ C$}
{$20^\circ C$}
\loigiai{$T(^\circ C)= T(K)- 273.15$, do đó, $300K = 26.85^\circ C \approx 27^\circ C$.}
\end{ex}

\dangbt{Chuyển đổi nhiệt độ giữa các thang đo thông dụng}
\begin{ex}
Một vật có nhiệt độ $373.15K$ sẽ tương đương với nhiệt độ nào trong thang Celsius?
\choice
{\True $100^\circ C$}
{$0^\circ C$}
{$273.15^\circ C$}
{$373.15^\circ C$}
\loigiai{$T(^\circ C)= T(K)- 273.15$, do đó, $373.15K = 100^\circ C$.}
\end{ex}

\dangbt{Chuyển đổi nhiệt độ giữa các thang đo thông dụng}
\begin{ex}
Nhiệt độ Celsius nào dưới đây tương ứng với tỉ số giữa nhiệt độ Fahrenheit và độ Celsius là $2,5$?
\choice
{$30^\circ C$}
{\True $45.71^\circ C$}
{$50^\circ C$}
{$60^\circ C$}
\loigiai{Ta đã giải bài toán và tìm ra nhiệt độ Celsius tương ứng với tỉ số giữa nhiệt độ Fahrenheit và độ Celsius là $2,5$. Từ phương trình: $F/C=2,5 \Rightarrow (32+1,8C)/C=2,5$ ta tìm được: $C \approx 45.71^\circ C$.}
\end{ex}

\dangbt{Chuyển đổi nhiệt độ giữa các thang đo thông dụng}
\begin{ex}
Nhiệt độ Kelvin nào dưới đây tương ứng với tỉ số giữa nhiệt độ Fahrenheit và độ Celsius là $3$?
\choice
{$273.15 K$}
{$298.15 K$}
{\True $299.82 K$}
{$310.15 K$}
\loigiai{Ta có $F/C=3$. Sử dụng công thức $F=9/5C+32$, giải phương trình $3C=9/5C+32 \Rightarrow C=26.67^\circ C$. Chuyển đổi sang Kelvin: $K=26.67+273.15=299.82K$.}
\end{ex}

\dangbt{Chuyển đổi nhiệt độ giữa các thang đo thông dụng}
\begin{ex}
Nhiệt độ Fahrenheit nào dưới đây tương ứng với tỉ số giữa nhiệt độ Fahrenheit và độ Celsius là $4$?
\choice
{$104^\circ F$}
{$128^\circ F$}
{\True $144^\circ F$}
{$160^\circ F$}
\loigiai{Ta có tỉ số $F/C =4$. Sử dụng công thức chuyển đổi từ Celsius sang Fahrenheit: $F=1.8C+32 \Rightarrow C \approx 14.55^\circ C \Rightarrow F=1.8C+32 \approx 144^\circ F$.}
\end{ex}

\dangbt{Chuyển đổi nhiệt độ giữa các thang đo thông dụng}
\begin{ex}
Nhiệt độ Fahrenheit nào dưới đây tương ứng với tỉ số giữa nhiệt độ Celsius và độ Kelvin là $0.5$?
\choice
{$100^\circ F$}
{$212^\circ F$}
{$300^\circ F$}
{\True $523.67^\circ F$}
\loigiai{Xác định tỉ số giữa độ Celsius (C) và độ Kelvin (K): $C/K=0.5$. Ta có $K=C+273.15 \Rightarrow C/(C+273.15)=0.5 \Rightarrow C=273.15^\circ C$. Chuyển đổi sang Fahrenheit: $F=1.8C+32 = 1.8 \times 273.15 + 32 = 491.67 + 32 = 523.67^\circ F$.}
\end{ex}

\dangbt{Chuyển đổi nhiệt độ giữa các thang đo thông dụng}
\begin{ex}
Nhiệt độ Fahrenheit của một vật giảm đi bao nhiêu độ khi từ $100^\circ C$ xuống $20^\circ C$?
\choice
{$68^\circ F$}
{$92^\circ F$}
{\True $144^\circ F$}
{$180^\circ F$}
\loigiai{Mức giảm độ F $\Delta F=1.8\Delta t = 1.8 \times (100-20) = 1.8 \times 80 = 144^\circ F$.}
\end{ex}

\dangbt{Chuyển đổi nhiệt độ giữa các thang đo thông dụng}
\begin{ex}
Nhiệt độ Fahrenheit nào dưới đây tương ứng với tỉ số giữa nhiệt độ Celsius và độ Kelvin là $0.25$?
\choice
{$100^\circ F$}
{$150^\circ F$}
{\True $196^\circ F$}
{$250^\circ F$}
\loigiai{Đầu tiên, ta biết các công thức chuyển đổi giữa các đơn vị nhiệt độ: $K=C+273.15$ và $F=1.8C+32$. Theo bài toán, tỉ số giữa nhiệt độ Celsius và độ Kelvin là $0.25$: $C/K=C/(C+273.15)=0.25 \Rightarrow C=0.25C+0.25 \times 273.15 \Rightarrow 0.75C = 68.2875 \Rightarrow C=91.05^\circ C$. Bây giờ, ta sẽ chuyển đổi nhiệt độ Celsius này sang Fahrenheit: $F=1.8C+32 = 1.8 \times 91.05 + 32 = 163.89 + 32 = 195.89^\circ F \approx 196^\circ F$.}
\end{ex}
